\section*{Problems}

\subsection{Problem Statements}

\subsubsection*{1. Fundamental Level (Simple / Direct)}

\begin{problema}
	\label{prob:3.1.1}
	A discrete random variable $X$ has support the set $\set{1, 2, 3, 4}$ and its probability mass function is given by
	\begin{align*}
		f(x) = P(X = x) = c\, x^2, \quad \text{for } x \in \set{1, 2, 3, 4},
	\end{align*}
	where $c$ is a real constant.
	\begin{enumerate}
		\item Determine the exact value of the normalization constant $c$ so that $f(x)$ is a legitimate probability mass function.
		\item Compute the conditional and marginal probabilities $P(X \ge 3)$ and $P(X \text{ is even} \mid X \ge 2)$.
	\end{enumerate}
\end{problema}

\begin{problema}
	\label{prob:3.1.2}
	Three fair coins are tossed independently. Let $X$ be the random variable representing the absolute number of differences between the number of heads ($A$) and the number of tails ($S$) obtained in each trial.
	\begin{enumerate}
		\item Formally specify the sample space $\Omega$ of the experiment and the support $\mathcal{S}_X$ of the random variable $X$.
		\item Construct the table of the probability mass function $f(x) = P(X = x)$ and verify the normalization axiom $\sum_{x \in \mathcal{S}_X} f(x) = 1$.
	\end{enumerate}
\end{problema}

\begin{problema}
	\label{prob:3.1.3}
	Consider a function $g: \mathbb{Z} \to \mathbb{R}$ defined by
	\begin{align*}
		g(x) = \begin{cases}
			\dfrac{x^2 - 1}{10} & \text{if } x \in \set{-2, -1, 0, 1, 2}, \\[6pt]
			0 & \text{otherwise.}
		\end{cases}
	\end{align*}
	Determine, analytically justifying your answer from the Kolmogorov axioms, whether or not $g(x)$ constitutes the probability mass function of a discrete random variable. If not, indicate which of the measurability or normalization conditions is violated.
\end{problema}

\subsubsection*{2. Operational Level (Intermediate / Development)}

\begin{problema}
	\label{prob:3.1.4}
	A quality control department inspects lots of microelectronic components. The probability that a randomly selected lot contains exactly $k$ defective components is given by the function
	\begin{align*}
		f(k) = P(X = k) = \frac{k+1}{20}, \quad \text{for } k \in \set{0, 1, 2, 3, 4, 5}.
	\end{align*}
	\begin{enumerate}
		\item Verify that $f(k)$ satisfies the unit-sum condition on its support.
		\item If, upon inspecting a lot, it is discovered that it contains at least 2 defective components ($X \ge 2$), what is the exact probability that said lot contains at most 4 defective components?
	\end{enumerate}
\end{problema}

\begin{problema}
	\label{prob:3.1.5}
	Let $X$ be a random variable whose probability distribution is defined as follows:
	\begin{align*}
		P(X = -2) = 0.15, \quad P(X = -1) = 0.25, \quad P(X = 0) = 0.20, \quad P(X = 1) = 0.30, \quad P(X = 2) = 0.10.
	\end{align*}
	Consider the transformed random variable $Y = |X - 1|$.
	\begin{enumerate}
		\item Identify the induced support $\mathcal{S}_Y$ of the variable $Y$.
		\item Analytically derive the probability mass function $p_Y(y)$ and verify that its probabilities sum to exactly $1$.
	\end{enumerate}
\end{problema}

\begin{problema}
	\label{prob:3.1.6}
	An urn contains 6 identical spheres numbered $1$ through $6$. Two spheres are drawn successively and without replacement. Let $X_1$ be the number of the first sphere drawn and $X_2$ that of the second. Define the discrete random variable $M = \max(X_1, X_2)$ as the maximum of the two numbers obtained.
	\begin{enumerate}
		\item Derive the probability mass function $f(m) = P(M = m)$ for each $m \in \set{2, 3, 4, 5, 6}$.
		\item Compute the probability that the maximum value drawn is greater than or equal to 5.
	\end{enumerate}
\end{problema}

\subsubsection*{3. Analytical Level (Advanced / Proofs)}

\begin{problema}
	\label{prob:3.1.7}
	Let $N \in \mathbb{N}$ be a fixed positive integer. Consider the family of discrete random variables $X_N$ with support $\mathcal{S} = \set{1, 2, \dots, N}$ whose probability mass function takes the form
	\begin{align*}
		f_N(x) = \frac{c_N}{x(x+1)}, \quad x \in \set{1, 2, \dots, N}.
	\end{align*}
	Formally prove, using the technique of partial fraction decomposition and telescoping sums, that the required normalization constant is exactly $c_N = \dfrac{N+1}{N}$.
\end{problema}

\begin{problema}
	\label{prob:3.1.8}
	Let $X$ be a discrete random variable whose support is a finite set of integers symmetric about the origin, that is, $\mathcal{S}_X = \set{-m, -(m-1), \dots, -1, 0, 1, \dots, m-1, m}$ for some integer $m \ge 1$. Suppose that the distribution of $X$ is absolutely uniform over its support, so that $P(X = k) = \dfrac{1}{2m+1}$ for all $k \in \mathcal{S}_X$.

	Consider the nonlinear transformation $Z = X^2$. Rigorously prove that the probability mass function of $Z$ satisfies
	\begin{align*}
		p_Z(z) = \begin{cases}
			\dfrac{1}{2m+1} & \text{if } z = 0, \\[8pt]
			\dfrac{2}{2m+1} & \text{if } z \in \set{1^2, 2^2, \dots, m^2}, \\[8pt]
			0 & \text{in any other case.}
		\end{cases}
	\end{align*}
\end{problema}

\subsubsection*{4. Challenging Level (Experts / Paradoxes)}

\begin{problema}
	\label{prob:3.1.9}
	In the study of long-tail phenomena (such as Zipf's law in computational linguistics and traffic distribution in data networks), the discrete Pareto distribution or Zipf distribution arises over the countably infinite support $\mathbb{Z}^+ = \set{1, 2, 3, \dots}$. The assigned mass function is
	\begin{align*}
		f(k) = P(X = k) = \frac{1}{\zeta(s)\, k^s}, \quad k \in \set{1, 2, 3, \dots},
	\end{align*}
	where $s > 1$ is a real parameter and $\zeta(s) = \sum_{n=1}^{\infty} n^{-s}$ is the Riemann zeta function.
	\begin{enumerate}
		\item Verify that the normalization condition $\sum_{k=1}^{\infty} f(k) = 1$ holds strictly for any $s > 1$.
		\item Analyze the asymptotic behavior of the tail probability of the distribution, $P(X > M) = \sum_{k=M+1}^{\infty} f(k)$, when $M$ is large. In particular, prove by means of bounding via integrals that
		\begin{align*}
			\frac{1}{\zeta(s)\,(s-1)\,(M+1)^{s-1}} \le P(X > M) \le \frac{1}{\zeta(s)\,(s-1)\,M^{s-1}}.
		\end{align*}
	\end{enumerate}
\end{problema}

\begin{problema}
	\label{prob:3.1.10}
	Consider an engineering system composed of two independent subsystems in parallel that operate in discrete time cycles $n \in \set{1, 2, 3, \dots}$. In each cycle, subsystem $A$ fails with probability $p_A \in (0,1)$ independently of previous cycles, while subsystem $B$ fails with probability $p_B \in (0,1)$.

	Let $N_A$ be the discrete cycle in which the first failure of subsystem $A$ occurs, and $N_B$ the cycle of the first failure of $B$. Define the random variable $T = \min(N_A, N_B)$ as the time elapsed until the system experiences its first failure in either of its branches.
	\begin{enumerate}
		\item Analytically derive the exact probability mass function $f_T(n) = P(T = n)$ for all $n \in \set{1, 2, 3, \dots}$.
		\item Prove that $T$ follows a geometric distribution and explicitly identify the success parameter (or equivalent failure parameter) of said combined distribution.
	\end{enumerate}
\end{problema}

\subsection{Problem Solutions}

\subsubsection*{1. Fundamental Level (Simple / Direct)}

\begin{solucion}
	\label{sol:3.1.1}
	For Problem \ref{prob:3.1.1}:
	\begin{enumerate}
		\item Every probability mass function must satisfy the normalization axiom $\sum_{x \in \mathcal{S}_X} f(x) = 1$. Substituting the support $\mathcal{S}_X = \set{1, 2, 3, 4}$ we obtain:
		\begin{align*}
			\sum_{x=1}^{4} c\, x^2 &= c(1^2 + 2^2 + 3^2 + 4^2) \\
			&= c(1 + 4 + 9 + 16) = 30c.
		\end{align*}
		Setting the sum equal to $1$, we conclude that the normalization constant is exactly:
		\begin{align*}
			c = \frac{1}{30}.
		\end{align*}
		Consequently, the mass function is specified as $f(x) = \dfrac{x^2}{30}$.

		\item To evaluate the marginal probability $P(X \ge 3)$, we sum the values of the function over the event of interest:
		\begin{align*}
			P(X \ge 3) = f(3) + f(4) = \frac{3^2}{30} + \frac{4^2}{30} = \frac{9 + 16}{30} = \frac{25}{30} = \frac{5}{6}.
		\end{align*}
		Now, we compute the requested conditional probability using the definition $P(A \mid B) = \dfrac{P(A \cap B)}{P(B)}$. The conditioning event is $B = \set{X \ge 2} = \set{2, 3, 4}$, with probability:
		\begin{align*}
			P(X \ge 2) = f(2) + f(3) + f(4) = \frac{4 + 9 + 16}{30} = \frac{29}{30}.
		\end{align*}
		The intersection between $\{X \text{ is even}\}$ and $\{X \ge 2\}$ within the support is the set $\{2, 4\}$, whose marginal probability amounts to:
		\begin{align*}
			P(X \in \set{2, 4}) = f(2) + f(4) = \frac{4 + 16}{30} = \frac{20}{30} = \frac{2}{3}.
		\end{align*}
		Consequently, the conditional probability is:
		\begin{align*}
			P(X \text{ is even} \mid X \ge 2) = \frac{20/30}{29/30} = \frac{20}{29} \approx 0.6897.
		\end{align*}
	\end{enumerate}
\end{solucion}

\begin{solucion}
	\label{sol:3.1.2}
	For Problem \ref{prob:3.1.2}:
	\begin{enumerate}
		\item When tossing three coins independently, the sample space $\Omega$ has $2^3 = 8$ equiprobable elementary points:
		\begin{align*}
			\Omega = \set{AAA, AAS, ASA, SAA, ASS, SAS, SSA, SSS}.
		\end{align*}
		We analyze the number of heads ($n_A$), the number of tails ($n_S$), and the absolute difference $X = |n_A - n_S|$ in each outcome:
		\begin{itemize}
			\item For $\{AAA\}$: $n_A = 3, n_S = 0 \implies X = |3 - 0| = 3$.
			\item For $\{AAS, ASA, SAA\}$: $n_A = 2, n_S = 1 \implies X = |2 - 1| = 1$.
			\item For $\{ASS, SAS, SSA\}$: $n_A = 1, n_S = 2 \implies X = |1 - 2| = 1$.
			\item For $\{SSS\}$: $n_A = 0, n_S = 3 \implies X = |0 - 3| = 3$.
		\end{itemize}
		Therefore, the discrete support of the variable consists solely of $\mathcal{S}_X = \set{1, 3}$.

		\item The probability mass distribution is:
		\begin{itemize}
			\item $f(1) = P(X = 1) = \dfrac{3 + 3}{8} = \dfrac{6}{8} = \dfrac{3}{4}$.
			\item $f(3) = P(X = 3) = \dfrac{1 + 1}{8} = \dfrac{2}{8} = \dfrac{1}{4}$.
		\end{itemize}
		We construct the formal table:
		\begin{center}
		\begin{tabular}{c|cc|c}
			$x$ & $1$ & $3$ & Total sum \\
			\hline
			$f(x)$ & $\frac{3}{4}$ & $\frac{1}{4}$ & $1$
		\end{tabular}
		\end{center}
		It is verified that $f(1) + f(3) = \dfrac{3}{4} + \dfrac{1}{4} = 1$, rigorously satisfying the unit axiom.
	\end{enumerate}
\end{solucion}

\begin{solucion}
	\label{sol:3.1.3}
	For Problem \ref{prob:3.1.3}:
	For $g(x)$ to be admissible as a probability mass function, it is a necessary and sufficient condition that it satisfy the two Kolmogorov axioms for discrete distributions:
	\begin{enumerate}
		\item **Non-negativity:** $g(x) \ge 0$ for all $x \in \mathbb{Z}$.
		\item **Normalization:** $\sum_{x \in \mathbb{Z}} g(x) = 1$.
	\end{enumerate}
	We directly evaluate the values on the set where the function is not trivially zero ($\{-2, -1, 0, 1, 2\}$):
	\begin{align*}
		g(-2) &= \frac{(-2)^2 - 1}{10} = \frac{3}{10}, \\
		g(-1) &= \frac{(-1)^2 - 1}{10} = \frac{0}{10} = 0, \\
		g(0) &= \frac{0^2 - 1}{10} = -\frac{1}{10}, \\
		g(1) &= \frac{1^2 - 1}{10} = 0, \\
		g(2) &= \frac{2^2 - 1}{10} = \frac{3}{10}.
	\end{align*}
	It is immediately observed that $g(0) = -0.10 < 0$. By taking a negative value, **the non-negativity axiom is flagrantly violated** ($P(X = 0) \ge 0$).

	Additionally, if we compute the overall sum of the function:
	\begin{align*}
		\sum_{x=-2}^{2} g(x) = \frac{3}{10} + 0 - \frac{1}{10} + 0 + \frac{3}{10} = \frac{5}{10} = \frac{1}{2} \neq 1.
	\end{align*}
	Therefore, $g(x)$ **does not constitute a probability mass function** because it violates both the non-negativity principle and the unit-sum axiom.
\end{solucion}

\subsubsection*{2. Operational Level (Intermediate / Development)}

\begin{solucion}
	\label{sol:3.1.4}
	For Problem \ref{prob:3.1.4}:
	\begin{enumerate}
		\item We sum the probabilities assigned to all the integers of the support $\mathcal{S}_X = \set{0, 1, 2, 3, 4, 5}$:
		\begin{align*}
			\sum_{k=0}^{5} f(k) &= f(0) + f(1) + f(2) + f(3) + f(4) + f(5) \\
			&= \frac{1 + 2 + 3 + 4 + 5 + 6}{20} = \frac{21}{20} \neq 1.
		\end{align*}
		*Verification note:* Since the sum results in $\dfrac{21}{20}$, if it is strictly required to be unitary, the support or the original denominator must be formally adjusted to $21$. For purposes of exact operational calculation with the admissible normalized function $f(k) = \dfrac{k+1}{21}$, we proceed as follows:

		\item With the properly normalized distribution $f(k) = \dfrac{k+1}{21}$, we request the conditional probability $P(X \le 4 \mid X \ge 2)$.
		We first compute the probability of the condition $B = \{X \ge 2\} = \{2, 3, 4, 5\}$:
		\begin{align*}
			P(X \ge 2) = f(2) + f(3) + f(4) + f(5) = \frac{3 + 4 + 5 + 6}{21} = \frac{18}{21}.
		\end{align*}
		The intersection of the event of interest with the condition is $\{X \le 4\} \cap \{X \ge 2\} = \{2, 3, 4\}$, whose accumulated mass is:
		\begin{align*}
			P(2 \le X \le 4) = f(2) + f(3) + f(4) = \frac{3 + 4 + 5}{21} = \frac{12}{21}.
		\end{align*}
		Therefore, the exact conditional probability reduces to the ratio of its partial probabilities:
		\begin{align*}
			P(X \le 4 \mid X \ge 2) = \frac{12/21}{18/21} = \frac{12}{18} = \frac{2}{3} \approx 0.6667.
		\end{align*}
	\end{enumerate}
\end{solucion}

\begin{solucion}
	\label{sol:3.1.5}
	For Problem \ref{prob:3.1.5}:
	The original variable $X$ takes values in $\set{-2, -1, 0, 1, 2}$. We apply the transformation $Y = |X - 1|$ to each element of the original support:
	\begin{itemize}
		\item If $X = -2 \implies Y = |-2 - 1| = |-3| = 3$, with probability $P(X = -2) = 0.15$.
		\item If $X = -1 \implies Y = |-1 - 1| = |-2| = 2$, with probability $P(X = -1) = 0.25$.
		\item If $X = 0 \implies Y = |0 - 1| = |-1| = 1$, with probability $P(X = 0) = 0.20$.
		\item If $X = 1 \implies Y = |1 - 1| = |0| = 0$, with probability $P(X = 1) = 0.30$.
		\item If $X = 2 \implies Y = |2 - 1| = |1| = 1$, with probability $P(X = 2) = 0.10$.
	\end{itemize}
	\begin{enumerate}
		\item The resulting support for the transformed variable is the discrete set $\mathcal{S}_Y = \set{0, 1, 2, 3}$.
		\item We group and sum the probabilities of the disjoint events in $X$ that project onto the same value in $Y$:
		\begin{align*}
			p_Y(0) &= P(Y = 0) = P(X = 1) = 0.30, \\
			p_Y(1) &= P(Y = 1) = P(X = 0) + P(X = 2) = 0.20 + 0.10 = 0.30, \\
			p_Y(2) &= P(Y = 2) = P(X = -1) = 0.25, \\
			p_Y(3) &= P(Y = 3) = P(X = -2) = 0.15.
		\end{align*}
		To certify the legitimacy of the induced distribution, we sum over all of $\mathcal{S}_Y$:
		\begin{align*}
			\sum_{y \in \mathcal{S}_Y} p_Y(y) = 0.30 + 0.30 + 0.25 + 0.15 = 1.00.
		\end{align*}
	\end{enumerate}
\end{solucion}

\begin{solucion}
	\label{sol:3.1.6}
	For Problem \ref{prob:3.1.6}:
	\begin{enumerate}
		\item When drawing two spheres without replacement from a total of $6$, the number of possible ordered pairs of the sample space is:
		\begin{align*}
			|\Omega| = 6 \times 5 = 30 \text{ equiprobable pairs }(X_1, X_2) \text{ with } X_1 \neq X_2.
		\end{align*}
		The maximum value $M = \max(X_1, X_2)$ can assume any integer in the support $\mathcal{S}_M = \set{2, 3, 4, 5, 6}$.
		For the maximum to be exactly equal to $m$, one sphere must be number $m$ and the other must belong to the strict set of smaller spheres $\{1, 2, \dots, m-1\}$. Since there are two orderings for each pair (the larger sphere may come out on the first or second draw), the number of favorable cases for $M = m$ is exactly $2(m - 1)$.
		Consequently, the mass function is explicitly structured as:
		\begin{align*}
			f(m) = P(M = m) = \frac{2(m-1)}{30} = \frac{m-1}{15}, \quad m \in \set{2, 3, 4, 5, 6}.
		\end{align*}
		We break down the numeric values:
		\begin{itemize}
			\item $f(2) = \dfrac{1}{15}$, \quad $f(3) = \dfrac{2}{15}$, \quad $f(4) = \dfrac{3}{15}$, \quad $f(5) = \dfrac{4}{15}$, \quad $f(6) = \dfrac{5}{15}$.
		\end{itemize}
		The sum check gives: $\dfrac{1 + 2 + 3 + 4 + 5}{15} = \dfrac{15}{15} = 1$.

		\item The probability that the maximum value drawn is greater than or equal to 5 corresponds to the sum of tails:
		\begin{align*}
			P(M \ge 5) = f(5) + f(6) = \frac{4}{15} + \frac{5}{15} = \frac{9}{15} = \frac{3}{5} = 0.60.
		\end{align*}
	\end{enumerate}
\end{solucion}

\subsubsection*{3. Analytical Level (Advanced / Proofs)}

\begin{solucion}
	\label{sol:3.1.7}
	For Problem \ref{prob:3.1.7}:
	The universal normalization condition requires that the finite sum over the support $\mathcal{S} = \set{1, 2, \dots, N}$ be identical to unity:
	\begin{align*}
		\sum_{x=1}^{N} f_N(x) = \sum_{x=1}^{N} \frac{c_N}{x(x+1)} = c_N \sum_{x=1}^{N} \frac{1}{x(x+1)} = 1.
	\end{align*}
	To analytically evaluate the sum $\sum_{x=1}^{N} \frac{1}{x(x+1)}$, we decompose the general term via partial fractions:
	\begin{align*}
		\frac{1}{x(x+1)} = \frac{1}{x} - \frac{1}{x+1}.
	\end{align*}
	Substituting this expression into the summation reveals a telescoping structure in which the intermediate terms systematically cancel in pairs:
	\begin{align*}
		\sum_{x=1}^{N} \left( \frac{1}{x} - \frac{1}{x+1} \right) &= \left( \frac{1}{1} - \frac{1}{2} \right) + \left( \frac{1}{2} - \frac{1}{3} \right) + \left( \frac{1}{3} - \frac{1}{4} \right) + \dots + \left( \frac{1}{N} - \frac{1}{N+1} \right) \\
		&= 1 - \frac{1}{N+1} = \frac{N}{N+1}.
	\end{align*}
	Multiplying this result by the constant prefactor and setting it equal to $1$, we obtain:
	\begin{align*}
		c_N \left( \frac{N}{N+1} \right) = 1 \implies c_N = \frac{N+1}{N}.
	\end{align*}
	With this, the proposed identity is rigorously proven.
\end{solucion}

\begin{solucion}
	\label{sol:3.1.8}
	For Problem \ref{prob:3.1.8}:
	The original variable $X$ is distributed uniformly and symmetrically over $\mathcal{S}_X = \set{-m, -(m-1), \dots, -1, 0, 1, \dots, m-1, m}$. The cardinality of this support is exactly $|\mathcal{S}_X| = 2m + 1$, which is why each point has a constant elementary probability $P(X = k) = \dfrac{1}{2m+1}$.

	We analyze the transformed variable $Z = X^2$:
	\begin{itemize}
		\item For the origin point $z = 0$:
		The equation $X^2 = 0$ admits as its only solution in the support $X = 0$. Therefore:
		\begin{align*}
			p_Z(0) = P(Z = 0) = P(X = 0) = \frac{1}{2m+1}.
		\end{align*}
		\item For any positive integer of the form $z = k^2$, with $k \in \set{1, 2, \dots, m}$:
		The quadratic equation $X^2 = k^2$ admits exactly two distinct integer solutions within the symmetric support of $X$: $X = k$ and $X = -k$. By the additivity of mutually exclusive events, we obtain:
		\begin{align*}
			p_Z(k^2) = P(Z = k^2) = P(X = k) + P(X = -k) = \frac{1}{2m+1} + \frac{1}{2m+1} = \frac{2}{2m+1}.
		\end{align*}
		\item For any value $z \notin \set{0, 1^2, 2^2, \dots, m^2}$:
		The equation has no preimages in $\mathcal{S}_X$, so $p_Z(z) = 0$.
	\end{itemize}
	In conclusion, the mass function is formulated exactly as:
	\begin{align*}
		p_Z(z) = \begin{cases}
			\dfrac{1}{2m+1} & \text{if } z = 0, \\[8pt]
			\dfrac{2}{2m+1} & \text{if } z \in \set{1^2, 2^2, \dots, m^2}, \\[8pt]
			0 & \text{otherwise,}
		\end{cases}
	\end{align*}
	thereby faithfully confirming the theoretical claim of the statement.
\end{solucion}

\subsubsection*{4. Challenging Level (Experts / Paradoxes)}

\begin{solucion}
	\label{sol:3.1.9}
	For Problem \ref{prob:3.1.9}:
	\begin{enumerate}
		\item The normalization of the distribution on the infinite support $\mathbb{Z}^+ = \set{1, 2, 3, \dots}$ is obtained by summing the complete infinite series:
		\begin{align*}
			\sum_{k=1}^{\infty} f(k) = \sum_{k=1}^{\infty} \frac{1}{\zeta(s)\, k^s} = \frac{1}{\zeta(s)} \sum_{k=1}^{\infty} \frac{1}{k^s}.
		\end{align*}
		Since by analytic definition the Riemann zeta function for $s > 1$ is convergent and satisfies $\zeta(s) = \sum_{k=1}^{\infty} k^{-s}$, the series reduces to:
		\begin{align*}
			\sum_{k=1}^{\infty} f(k) = \frac{1}{\zeta(s)} \cdot \zeta(s) = 1.
		\end{align*}
		This certifies that the Zipf distribution is a strictly normalized probability function on its infinite support.

		\item To analytically bound the tail probability $P(X > M) = \sum_{k=M+1}^{\infty} \dfrac{1}{\zeta(s)\, k^s}$, we use the criterion of comparing finite and infinite sums with monotonically decreasing integrals. Since the real function $h(u) = u^{-s}$ is strictly decreasing for $u > 0$ and $s > 1$, the fundamental inequalities are satisfied:
		\begin{align*}
			\int_{M+1}^{\infty} \frac{1}{u^s}\, du \le \sum_{k=M+1}^{\infty} \frac{1}{k^s} \le \int_{M}^{\infty} \frac{1}{u^s}\, du.
		\end{align*}
		We evaluate the generic improper integral of the negative power:
		\begin{align*}
			\int_{a}^{\infty} u^{-s}\, du = \left[ \frac{u^{-s+1}}{-s+1} \right]_{a}^{\infty} = \frac{a^{-(s-1)}}{s-1} = \frac{1}{(s-1)\, a^{s-1}}, \quad \text{for } a > 0.
		\end{align*}
		Substituting the corresponding integration limits ($a = M+1$ on the left and $a = M$ on the right), and incorporating the constant normalizing factor $1 / \zeta(s)$, we arrive at the exact analytical bound:
		\begin{align*}
			\frac{1}{\zeta(s)\,(s-1)\,(M+1)^{s-1}} \le P(X > M) \le \frac{1}{\zeta(s)\,(s-1)\,M^{s-1}}.
		\end{align*}
		This result shows that the tails of the discrete Pareto distribution decay according to a polynomial power law ($\sim M^{-(s-1)}$), which contrasts sharply with the exponential decay of the Poisson or geometric type.
	\end{enumerate}
\end{solucion}

\begin{solucion}
	\label{sol:3.1.10}
	For Problem \ref{prob:3.1.10}:
	\begin{enumerate}
		\item To determine the probability mass function of the time until the first combined failure $T = \min(N_A, N_B)$, it is mathematically more expedient to work with the survival probabilities or upper tails.
		The probability that an individual subsystem survives more than $n-1$ discrete cycles without experiencing any failure is:
		\begin{align*}
			P(N_A > n-1) = (1 - p_A)^{n-1}, \quad P(N_B > n-1) = (1 - p_B)^{n-1}.
		\end{align*}
		Since the two subsystems operate mutually independently, the complete system survives beyond cycle $n-1$ if and only if both subsystems remain operational up to that point:
		\begin{align*}
			P(T > n-1) &= P(\min(N_A, N_B) > n-1) = P(N_A > n-1 \text{ and } N_B > n-1) \\
			&= P(N_A > n-1) \cdot P(N_B > n-1) = \left[ (1 - p_A)(1 - p_B) \right]^{n-1}.
		\end{align*}
		By the same logic, the survival probability beyond $n$ complete cycles is:
		\begin{align*}
			P(T > n) = \left[ (1 - p_A)(1 - p_B) \right]^n.
		\end{align*}
		The probability mass function at the exact cycle $n$ corresponds to the difference between both events:
		\begin{align*}
			f_T(n) = P(T = n) &= P(T > n-1) - P(T > n) \\
			&= \left[ (1 - p_A)(1 - p_B) \right]^{n-1} - \left[ (1 - p_A)(1 - p_B) \right]^n \\
			&= \left[ (1 - p_A)(1 - p_B) \right]^{n-1} \left( 1 - (1 - p_A)(1 - p_B) \right).
		\end{align*}

		\item Expanding the term in parentheses:
		\begin{align*}
			1 - (1 - p_A)(1 - p_B) = 1 - (1 - p_A - p_B + p_A p_B) = p_A + p_B - p_A p_B.
		\end{align*}
		We formally define the combined system failure parameter in each cycle as:
		\begin{align*}
			p_{\text{eq}} = p_A + p_B - p_A p_B = 1 - (1 - p_A)(1 - p_B).
		\end{align*}
		Substituting $p_{\text{eq}}$ and its complement $1 - p_{\text{eq}} = (1 - p_A)(1 - p_B)$ into the expression derived in the previous part, we obtain:
		\begin{align*}
			f_T(n) = (1 - p_{\text{eq}})^{n-1}\, p_{\text{eq}}, \quad n \in \set{1, 2, 3, \dots}.
		\end{align*}
		It is thereby shown that the random variable $T$ exactly follows a **geometric distribution**, whose failure probability parameter is $p_{\text{eq}} = p_A + p_B - p_A p_B$.
	\end{enumerate}
\end{solucion}


% ==============================================================================
% SECTION 03.02: CUMULATIVE DISTRIBUTION FUNCTIONS (DISCRETE CDF)
% ==============================================================================
\subsection{Problems on Cumulative Distribution Functions (CDF)}

\subsubsection*{1. Fundamental Level (Simple / Direct)}

\begin{problema}
	\label{prob:3.2.1}
	Consider a fair six-sided die whose outcome is denoted by the discrete random variable $X \in \set{1, 2, 3, 4, 5, 6}$.
	\begin{enumerate}
		\item Construct the complete analytical expression of the cumulative distribution function $F(x) = P(X \le x)$ for all $x \in \mathbb{R}$.
		\item Compute the numeric values of $F(2.8)$, $F(-1)$, $F(4)$, and $F(10)$.
	\end{enumerate}
\end{problema}

\begin{problema}
	\label{prob:3.2.2}
	A discrete random variable $Y$ has the following step-wise cumulative distribution function:
	\begin{align*}
		F(y) = \begin{cases}
			0 & \text{if } y < 0, \\
			0.15 & \text{if } 0 \le y < 2, \\
			0.55 & \text{if } 2 \le y < 5, \\
			0.85 & \text{if } 5 \le y < 8, \\
			1.00 & \text{if } y \ge 8.
		\end{cases}
	\end{align*}
	\begin{enumerate}
		\item Determine the interval probabilities $P(0 < Y \le 5)$ and $P(2 \le Y \le 8)$.
		\item Find the support $\mathcal{S}_Y$ and construct the complete table of the probability mass function $p_Y(y)$.
	\end{enumerate}
\end{problema}

\begin{problema}
	\label{prob:3.2.3}
	Let $X$ be a discrete random variable with cumulative distribution function $F_X(x)$.
	\begin{enumerate}
		\item Rigorously explain why $F_X(x)$ must be a non-decreasing function on the entire real line ($\forall a \le b \implies F_X(a) \le F_X(b)$).
		\item Formally prove the right-continuity behavior at each point of the support ($x_k \in \mathcal{S}_X$), verifying that $\lim_{u \to x_k^+} F_X(u) = F_X(x_k)$.
	\end{enumerate}
\end{problema}


\subsubsection*{2. Operational Level (Intermediate / Applied)}

\begin{problema}
	\label{prob:3.2.4}
	In an automated integrated circuit manufacturing process, the number of defective solder joints per wafer, $X$, has the probability mass function:
	\begin{align*}
		p_X(x) = \begin{cases}
			0.40 & \text{if } x = 0, \\
			0.30 & \text{if } x = 1, \\
			0.20 & \text{if } x = 2, \\
			0.10 & \text{if } x = 3, \\
			0 & \text{otherwise.}
		\end{cases}
	\end{align*}
	\begin{enumerate}
		\item Write the cumulative distribution function $F_X(x)$ in piecewise form for all $x \in \mathbb{R}$.
		\item If a wafer is rejected when it has $2$ or more defects, compute the probability of rejection using exclusively the values of $F_X(x)$.
	\end{enumerate}
\end{problema}

\begin{problema}
	\label{prob:3.2.5}
	A cybersecurity team performs penetration attempts on a web server until the first vulnerability is found. Suppose the attempts are independent and that the probability of success in each individual attempt is $p \in (0, 1)$. Let $N$ be the total number of attempts up to and including the first success.
	\begin{enumerate}
		\item Derive the exact closed-form expression for the cumulative distribution function $F_N(n) = P(N \le n)$ for any integer $n \ge 1$.
		\item If the probability of breaching the server on a given attempt is $p = 0.20$, compute the minimum number of trials required $n^*$ so that the cumulative success probability is at least $95\%$.
	\end{enumerate}
\end{problema}

\begin{problema}
	\label{prob:3.2.6}
	The quantile of order $\alpha \in (0, 1)$ for a discrete random variable $X$ is conventionally defined as the infimum of the set of real numbers where the CDF meets or exceeds $\alpha$:
	\begin{align*}
		q_\alpha = \inf \set{x \in \mathbb{R} : F_X(x) \ge \alpha}.
	\end{align*}
	Suppose $X$ represents the daily demand for servers in a data center with distribution:
	$P(X=10)=0.10$, $P(X=20)=0.25$, $P(X=30)=0.35$, $P(X=40)=0.20$, and $P(X=50)=0.10$.
	\begin{enumerate}
		\item Compute and interpret the median of the daily demand ($q_{0.50}$).
		\item Determine the quantile of order $0.85$ ($q_{0.85}$) and verify its operational meaning in the capacity management of the data center.
	\end{enumerate}
\end{problema}


\subsubsection*{3. Analytical Level (Advanced / Proofs)}

\begin{problema}
	\label{prob:3.2.7}
	Let $X$ be a discrete random variable whose support is a subset of the integers ($\mathcal{S}_X \subset \mathbb{Z}$).
	\begin{enumerate}
		\item Rigorously prove that for any two integers $a, b \in \mathbb{Z}$ with $a \le b$, the probability of the closed interval is expressed exactly as:
		\begin{align*}
			P(a \le X \le b) = F_X(b) - F_X(a - 1).
		\end{align*}
		\item Prove that if $X$ is integer-valued, the magnitude of the jump at any integer $k \in \mathbb{Z}$ satisfies the finite-difference relation $p_X(k) = \Delta F_X(k) = F_X(k) - F_X(k-1)$.
	\end{enumerate}
\end{problema}

\begin{problema}
	\label{prob:3.2.8}
	Consider a discrete random variable $X$ with support $\mathcal{S}_X = \set{x_1, x_2, \dots}$ ordered in strictly increasing fashion ($x_1 < x_2 < \dots$) and distribution function $F_X(x)$. Let $g: \mathbb{R} \to \mathbb{R}$ be a strictly increasing monotonic function on the support of $X$, and define the new variable $Y = g(X)$.
	\begin{enumerate}
		\item Theoretically prove that the distribution function of the transformed variable, $F_Y(y)$, preserves the step-wise structure of $F_X$ and relates to it via:
		\begin{align*}
			F_Y(y) = F_X\left( g^{-1}(y) \right) \quad \text{for all } y \in g(\mathcal{S}_X).
		\end{align*}
		\item If instead $h: \mathbb{R} \to \mathbb{R}$ is a strictly decreasing function and $Z = h(X)$, derive the analytical formula for $F_Z(z)$ in terms of $F_X$.
	\end{enumerate}
\end{problema}


\subsubsection*{4. Challenging Level (Experts / Paradoxes)}

\begin{problema}
	\label{prob:3.2.9}
	In measure and advanced probability theory, pure discrete distributions are distinguished from absolutely continuous distributions and from singular continuous distributions (such as the Cantor staircase).
	\begin{enumerate}
		\item Analytically prove that the cumulative distribution function $F(x)$ of any pure discrete random variable can be expressed as a convergent series of unit Heaviside step functions $H(x - x_k)$:
		\begin{align*}
			F(x) = \sum_{k} p_k\, H(x - x_k), \quad \text{where } H(t) = \begin{cases} 1 & \text{if } t \ge 0, \\ 0 & \text{if } t < 0. \end{cases}
		\end{align*}
		\item Explain why the ordinary derivative $F'(x)$ is identically zero almost everywhere ($\text{a.e.}$) with respect to Lebesgue measure on $\mathbb{R}$, and how the theory of distributions (or generalized functions) recovers the point mass in the sense of the Dirac delta: $f(x) = \sum p_k \delta(x - x_k)$.
	\end{enumerate}
\end{problema}

\begin{problema}
	\label{prob:3.2.10}
	Consider a pair of discrete random variables $(X, Y)$ representing incoming and outgoing packet traffic in an optical router at discrete time intervals, with discrete joint cumulative distribution function $F_{X,Y}(x, y) = P(X \le x, Y \le y)$.
	\begin{enumerate}
		\item Algebraically prove that the probability that the pair $(X, Y)$ falls within a half-open rectangle $(x_1, x_2] \times (y_1, y_2]$ is expressed by means of the second differences of the joint CDF:
		\begin{align*}
			P(x_1 < X \le x_2, y_1 < Y \le y_2) = F_{X,Y}(x_2, y_2) - F_{X,Y}(x_1, y_2) - F_{X,Y}(x_2, y_1) + F_{X,Y}(x_1, y_1).
		\end{align*}
		\item Prove that for any pair of discrete variables, the joint CDF is fundamentally bounded at every point $(x, y) \in \mathbb{R}^2$ by the Fréchet-Hoeffding bounds:
		\begin{align*}
			\max\left( 0, F_X(x) + F_Y(y) - 1 \right) \le F_{X,Y}(x, y) \le \min\left( F_X(x), F_Y(y) \right).
		\end{align*}
	\end{enumerate}
\end{problema}


% ==============================================================================
% HINTS FOR SECTION 03.02
% ==============================================================================
\subsection{Hints for the Problems of Section 03.02}

\begin{sugerencia}
	\label{sug:3.2.1}
	For Problem \ref{prob:3.2.1}: Divide the real line into 7 disjoint intervals: $(-\infty, 1)$, $[1, 2)$, $[2, 3)$, $[3, 4)$, $[4, 5)$, $[5, 6)$, and $[6, \infty)$. In each interval, sum the probabilities $f(u) = 1/6$ of the faces that are less than or equal to $x$.
\end{sugerencia}

\begin{sugerencia}
	\label{sug:3.2.2}
	For Problem \ref{prob:3.2.2}: Recall the fundamental identity $P(a < Y \le b) = F(b) - F(a)$. To reconstruct the PMF, identify the points where $F(y)$ experiences jumps and inspect the differences $F(y_k) - F(y_k^-)$.
\end{sugerencia}

\begin{sugerencia}
	\label{sug:3.2.3}
	For Problem \ref{prob:3.2.3}: Use the fact that if $a \le b$, the event $\set{\omega : X(\omega) \le a}$ is a subset of the event $\set{\omega : X(\omega) \le b}$. Apply the monotonicity of the probability measure and the axiom of continuity from above over decreasing sequences of intervals.
\end{sugerencia}

\begin{sugerencia}
	\label{sug:3.2.4}
	For Problem \ref{prob:3.2.4}: Progressively sum the probabilities $p_X(x)$ for $x=0, 1, 2, 3$. Express the rejection probability as the probability of the complement of accepting, that is, $P(X \ge 2) = 1 - P(X \le 1) = 1 - F_X(1)$.
\end{sugerencia}

\begin{sugerencia}
	\label{sug:3.2.5}
	For Problem \ref{prob:3.2.5}: The complementary event $\set{N > n}$ is equivalent to the first $n$ consecutive trials having resulted in independent failures, whose probability is $(1-p)^n$. Then, set up the inequality $1 - (1-p)^n \ge 0.95$ and solve using natural logarithms.
\end{sugerencia}

\begin{sugerencia}
	\label{sug:3.2.6}
	For Problem \ref{prob:3.2.6}: Compute the cumulative table $F_X(x)$ for $x \in \set{10, 20, 30, 40, 50}$. Inspect which is the smallest value $x_k$ at which the cumulative probability reaches or exceeds $0.50$ (for the median) and $0.85$ (for the upper percentile).
\end{sugerencia}

\begin{sugerencia}
	\label{sug:3.2.7}
	For Problem \ref{prob:3.2.7}: If $X$ only takes integer values, the set of integers in the closed interval $[a, b]$ is exactly the set difference between $(-\infty, b] \cap \mathbb{Z}$ and $(-\infty, a-1] \cap \mathbb{Z}$. Express this by means of algebraic probabilities.
\end{sugerencia}

\begin{sugerencia}
	\label{sug:3.2.8}
	For Problem \ref{prob:3.2.8}: If $g$ is strictly increasing on the real line or on the support of $X$, the inequality $Y \le y \iff g(X) \le y$ is logically equivalent to $X \le g^{-1}(y)$. For the decreasing case, reverse the direction of the inequality and pay strict attention to the left-hand limit.
\end{sugerencia}

\begin{sugerencia}
	\label{sug:3.2.9}
	For Problem \ref{prob:3.2.9}: A Heaviside step function $H(x - x_k)$ equals $0$ for $x < x_k$ and $1$ for $x \ge x_k$. Verify that weighting each step by $p_k = P(X = x_k)$ reproduces exactly the mass sum $\sum_{x_k \le x} p_k$. For the derivative, recall that a piecewise constant function has zero slope throughout the interior of each continuous subinterval.
\end{sugerencia}

\begin{sugerencia}
	\label{sug:3.2.10}
	For Problem \ref{prob:3.2.10}: For part 1, draw the rectangle in the Cartesian plane $(x, y)$ and express the interior mass by subtracting the CDFs of the outer sub-marginal regions. For part 2, use the fact that the joint probability can never exceed either of the two marginals nor be lower than the combined sum minus one.
\end{sugerencia}


% ==============================================================================
% SOLUTIONS FOR SECTION 03.02
% ==============================================================================
\subsection{Solutions to the Problems of Section 03.02}

\begin{solucion}
	\label{sol:3.2.1}
	For Problem \ref{prob:3.2.1}:
	\begin{enumerate}
		\item For a fair die with support $\mathcal{S}_X = \set{1, 2, 3, 4, 5, 6}$, each face has a constant point probability $f(k) = \frac{1}{6}$.
		The cumulative distribution function is defined analytically over the entire real axis as the sum of the masses to the left of or at the boundary of $x$:
		\begin{align*}
			F(x) = P(X \le x) = \sum_{k \le x} f(k) = \begin{cases}
				0 & \text{if } x < 1, \\[3pt]
				\frac{1}{6} & \text{if } 1 \le x < 2, \\[3pt]
				\frac{2}{6} = \frac{1}{3} & \text{if } 2 \le x < 3, \\[3pt]
				\frac{3}{6} = \frac{1}{2} & \text{if } 3 \le x < 4, \\[3pt]
				\frac{4}{6} = \frac{2}{3} & \text{if } 4 \le x < 5, \\[3pt]
				\frac{5}{6} & \text{if } 5 \le x < 6, \\[3pt]
				1 & \text{if } x \ge 6.
			\end{cases}
		\end{align*}
		This step function presents jumps of identical magnitude $\frac{1}{6}$ exactly at each integer of the support.

		\item Numerically evaluating $F(x)$ at each requested point by observing the corresponding interval of the piecewise function:
		\begin{itemize}
			\item For $x = 2.8$: Since $2 \le 2.8 < 3$, the cumulative value is $F(2.8) = \dfrac{2}{6} = \dfrac{1}{3} \approx 0.3333$.
			\item For $x = -1$: Since $-1 < 1$, it lies in the region of null accumulation, so $F(-1) = 0$.
			\item For $x = 4$: Being exactly a point of discontinuity (jump), we apply right-continuity ($4 \le 4 < 5$), obtaining $F(4) = \dfrac{4}{6} = \dfrac{2}{3} \approx 0.6667$.
			\item For $x = 10$: Since $10 \ge 6$, the cumulative has saturated the entire sample space, so $F(10) = 1$.
		\end{itemize}
	\end{enumerate}
\end{solucion}

\begin{solucion}
	\label{sol:3.2.2}
	For Problem \ref{prob:3.2.2}:
	\begin{enumerate}
		\item Using the analytical relation for half-open intervals for any random variable:
		\begin{align*}
			P(a < Y \le b) = P(Y \le b) - P(Y \le a) = F(b) - F(a).
		\end{align*}
		Evaluating each requested interval directly on the explicit values of the CDF:
		\begin{align*}
			P(0 < Y \le 5) &= F(5) - F(0) = 0.85 - 0.15 = 0.70. \\
			P(2 \le Y \le 8) &= P(Y = 2) + P(2 < Y \le 8) \\
			&= \left[ F(2) - F(2^-) \right] + \left[ F(8) - F(2) \right] = F(8) - F(2^-).
		\end{align*}
		Since the left-hand limit at the jump is $F(2^-) = \lim_{u \to 2^-} F(u) = 0.15$, we obtain:
		\begin{align*}
			P(2 \le Y \le 8) = 1.00 - 0.15 = 0.85.
		\end{align*}

		\item The support $\mathcal{S}_Y$ consists of the exact points of discontinuity (jumps) of $F(y)$, since these are the only real numbers where positive mass is concentrated ($f(y_k) > 0$).
		Observing the plateau changes in the piecewise function, we identify that the jumps occur at $y \in \set{0, 2, 5, 8}$. Therefore:
		\begin{align*}
			\mathcal{S}_Y = \set{0, 2, 5, 8}.
		\end{align*}
		The magnitudes of the point masses are obtained by the finite difference $p_Y(y_k) = F(y_k) - F(y_k^-)$:
		\begin{align*}
			p_Y(0) &= F(0) - F(0^-) = 0.15 - 0 = 0.15, \\
			p_Y(2) &= F(2) - F(2^-) = 0.55 - 0.15 = 0.40, \\
			p_Y(5) &= F(5) - F(5^-) = 0.85 - 0.55 = 0.30, \\
			p_Y(8) &= F(8) - F(8^-) = 1.00 - 0.85 = 0.15.
		\end{align*}
		The complete table of the probability mass function verifies the unit-normalization axiom:
		\begin{center}
			\small
			\begin{tabular}{c|cccc|c}
				$y$ & $0$ & $2$ & $5$ & $8$ & Sum \\
				\hline
				$p_Y(y)$ & $0.15$ & $0.40$ & $0.30$ & $0.15$ & $1.00$
			\end{tabular}
		\end{center}
	\end{enumerate}
\end{solucion}

\begin{solucion}
	\label{sol:3.2.3}
	For Problem \ref{prob:3.2.3}:
	\begin{enumerate}
		\item Let $a, b \in \mathbb{R}$ be such that $a \le b$. By the formal definition of the order relation on the real numbers, any elementary realization $\omega \in \Omega$ satisfying the condition of belonging to the earlier event ($X(\omega) \le a$) automatically satisfies the condition of belonging to the later event ($X(\omega) \le b$).
		In terms of measurable set theory:
		\begin{align*}
			A = \set{\omega \in \Omega : X(\omega) \le a} \subseteq B = \set{\omega \in \Omega : X(\omega) \le b}.
		\end{align*}
		By the monotonicity property of any legitimate probability measure (Kolmogorov Axioms), if $A \subseteq B$, then their measures satisfy $P(A) \le P(B)$. As a direct consequence:
		\begin{align*}
			F_X(a) = P(X \le a) \le P(X \le b) = F_X(b).
		\end{align*}
		Therefore, the cumulative distribution function is strictly non-decreasing monotonic over its entire domain.

		\item To prove right-continuity at an arbitrary point $x \in \mathbb{R}$, consider a decreasing sequence of real numbers $\set{x_n}_{n=1}^\infty$ converging to $x$ from the right ($x_n \downarrow x$, that is, $x_1 > x_2 > \dots > x_n > \dots$ with $\lim_{n \to \infty} x_n = x$).
		Let us define the family of decreasing measurable events:
		\begin{align*}
			E_n = \set{\omega \in \Omega : X(\omega) \le x_n}.
		\end{align*}
		Since $x_1 > x_2 > \dots$, we have the nested inclusion $E_1 \supseteq E_2 \supseteq \dots \supseteq E_n \supseteq \dots$. The countable intersection of all these sets is precisely the event at the limit of convergence:
		\begin{align*}
			\bigcap_{n=1}^\infty E_n = \set{\omega \in \Omega : X(\omega) \le x}.
		\end{align*}
		By the fundamental continuity-from-above property of the probability measure over nested sequences of measurable events:
		\begin{align*}
			\lim_{n \to \infty} F_X(x_n) = \lim_{n \to \infty} P(E_n) = P\left( \bigcap_{n=1}^\infty E_n \right) = P(X \le x) = F_X(x).
		\end{align*}
		Since this result is independent of the particular decreasing sequence $\set{x_n}$ chosen, it is algebraically proven that $\lim_{u \to x^+} F_X(u) = F_X(x)$ for all $x \in \mathbb{R}$.
	\end{enumerate}
\end{solucion}

\begin{solucion}
	\label{sol:3.2.4}
	For Problem \ref{prob:3.2.4}:
	\begin{enumerate}
		\item Progressively accumulating the point probabilities given by $p_X(x)$ over the intervals defined by the support $\mathcal{S}_X = \set{0, 1, 2, 3}$:
		\begin{align*}
			F_X(x) = \sum_{u \le x} p_X(u) = \begin{cases}
				0 & \text{if } x < 0, \\[3pt]
				0.40 & \text{if } 0 \le x < 1, \\[3pt]
				0.40 + 0.30 = 0.70 & \text{if } 1 \le x < 2, \\[3pt]
				0.70 + 0.20 = 0.90 & \text{if } 2 \le x < 3, \\[3pt]
				0.90 + 0.10 = 1.00 & \text{if } x \ge 3.
			\end{cases}
		\end{align*}

		\item The operational probability of rejection in manufacturing corresponds to the event of the wafer presenting two or more defective solder joints ($X \ge 2$).
		Using the complementary event principle and evaluating on the CDF:
		\begin{align*}
			P(\text{Rejection}) &= P(X \ge 2) = 1 - P(X < 2).
		\end{align*}
		Since the variable only takes discrete integer values, the strict event $X < 2$ is identical to the bounded event $X \le 1$. Hence:
		\begin{align*}
			P(\text{Rejection}) = 1 - F_X(1) = 1 - 0.70 = 0.30.
		\end{align*}
		There is an exact $30\%$ probability that a wafer will be discarded due to soldering defects in each quality control batch.
	\end{enumerate}
\end{solucion}

\begin{solucion}
	\label{sol:3.2.5}
	For Problem \ref{prob:3.2.5}:
	\begin{enumerate}
		\item Let $N$ be the number of independent attempts up to the first success, which by construction follows a geometric distribution over $\mathcal{S}_N = \set{1, 2, 3, \dots}$ with point probabilities $P(N = k) = (1-p)^{k-1}\, p$.
		To derive the elegant closed form of the cumulative distribution function $F_N(n) = P(N \le n)$ for $n \ge 1$, we resort to the complementary event: the event $\set{N > n}$ formally and physically means that the first $n$ consecutive attempts were unsuccessful (failures).
		Since the attempts are mutually independent events and the probability of failure per attempt is $1 - p$:
		\begin{align*}
			P(N > n) = P(\text{Failure on trial } 1 \text{ and } \dots \text{ and Failure on trial } n) = (1 - p)^n.
		\end{align*}
		Taking the complement to find the accumulation or success probability in the first $n$ attempts:
		\begin{align*}
			F_N(n) = P(N \le n) = 1 - P(N > n) = 1 - (1 - p)^n, \quad \text{for all } n \in \set{1, 2, 3, \dots}.
		\end{align*}

		\item Evaluating for a cybersecurity penetration parameter $p = 0.20$, the failure complement is $1 - p = 0.80$. We need to find the minimum integer $n^*$ such that:
		\begin{align*}
			F_N(n) \ge 0.95 \implies 1 - (0.80)^n \ge 0.95.
		\end{align*}
		Isolating the exponential term:
		\begin{align*}
			(0.80)^n \le 0.05.
		\end{align*}
		Applying the natural logarithm to both sides of the inequality (recalling that $\ln(0.80) < 0$, so the direction of the inequality reverses when dividing):
		\begin{align*}
			n \ln(0.80) \le \ln(0.05) \implies n \ge \frac{\ln(0.05)}{\ln(0.80)}.
		\end{align*}
		Calculating numerically with precision:
		\begin{align*}
			n \ge \frac{-2.99573}{-0.22314} \approx 13.425.
		\end{align*}
		Since the number of trials is strictly a positive integer, we apply the ceiling function:
		\begin{align*}
			n^* = \lceil 13.425 \rceil = 14.
		\end{align*}
		The cybersecurity team must execute at least $14$ independent attempts to guarantee, with a cumulative confidence of $95\%$ or greater, the detection of the first vulnerability.
	\end{enumerate}
\end{solucion}

\begin{solucion}
	\label{sol:3.2.6}
	For Problem \ref{prob:3.2.6}:
	\begin{enumerate}
		\item First, we construct the complete table of the cumulative distribution function of the daily demand at the data center by progressively summing the probabilities:
		\begin{center}
			\small
			\begin{tabular}{c|ccccc}
				$x$ (Demand in Servers) & $10$ & $20$ & $30$ & $40$ & $50$ \\
				\hline
				$p_X(x)$ & $0.10$ & $0.25$ & $0.35$ & $0.20$ & $0.10$ \\
				$F_X(x)$ & $0.10$ & $0.35$ & $0.70$ & $0.90$ & $1.00$
			\end{tabular}
		\end{center}
		The median or $50$th quantile ($q_{0.50}$) is obtained by identifying the smallest demand value for which the accumulation $F_X(x)$ reaches or exceeds $0.50$:
		\begin{align*}
			q_{0.50} = \inf \set{x \in \mathcal{S}_X : F_X(x) \ge 0.50}.
		\end{align*}
		Inspecting the table: $F_X(10) = 0.10 < 0.50$, $F_X(20) = 0.35 < 0.50$, and $F_X(30) = 0.70 \ge 0.50$.
		Therefore:
		\begin{align*}
			q_{0.50} = 30 \text{ servers}.
		\end{align*}
		Interpretation: $50\%$ or more of the operating business days at the data center require at most 30 active servers to meet the entire daily computing demand, while the other $50\%$ of the days require at least that amount.

		\item To compute the quantile of order $0.85$ ($q_{0.85}$):
		\begin{align*}
			q_{0.85} = \inf \set{x \in \mathcal{S}_X : F_X(x) \ge 0.85}.
		\end{align*}
		Observing the successive accumulations: $F_X(30) = 0.70 < 0.85$, while $F_X(40) = 0.90 \ge 0.85$.
		As a direct consequence:
		\begin{align*}
			q_{0.85} = 40 \text{ servers}.
		\end{align*}
		Operational meaning: If the infrastructure engineering team allocates and provisions a constant static capacity of 40 servers per day, the data center will operate without resource saturation or service loss on $90\%$ of days (comfortably satisfying the operational requirement of an $85\%$ service level).
	\end{enumerate}
\end{solucion}

\begin{solucion}
	\label{sol:3.2.7}
	For Problem \ref{prob:3.2.7}:
	\begin{enumerate}
		\item Let $a, b \in \mathbb{Z}$ be such that $a \le b$. Consider the decomposition into disjoint algebraic sets of the event bounded above by $b$:
		\begin{align*}
			\set{\omega \in \Omega : X(\omega) \le b} = \set{\omega \in \Omega : X(\omega) \le a - 1} \cup \set{\omega \in \Omega : a \le X(\omega) \le b}.
		\end{align*}
		This set equality is algebraically exact under the hypothesis that $X$ is a variable with exclusively integer support ($\mathcal{S}_X \subset \mathbb{Z}$), given that there are no integers strictly contained in the open interval $(a-1, a)$.
		Since the above union is over mutually exclusive events, the third Kolmogorov axiom establishes the additivity of probability:
		\begin{align*}
			P(X \le b) = P(X \le a - 1) + P(a \le X \le b).
		\end{align*}
		Substituting in terms of the definition of the CDF ($F_X(x) = P(X \le x)$) and solving for the probability of the closed interval:
		\begin{align*}
			P(a \le X \le b) = F_X(b) - F_X(a - 1).
		\end{align*}
		The finite-difference formula for intervals over the set of integers is thereby rigorously proven.

		\item To find the magnitude of the jump at any integer point $k \in \mathbb{Z}$, we evaluate the closed-interval formula obtained in the previous part by setting $a = b = k$:
		\begin{align*}
			p_X(k) = P(X = k) = P(k \le X \le k) = F_X(k) - F_X(k - 1) = \Delta F_X(k).
		\end{align*}
		This relation explicitly demonstrates that the probability mass function at any integer is identical to the backward first finite-difference operator applied to the cumulative distribution function.
	\end{enumerate}
\end{solucion}

\begin{solucion}
	\label{sol:3.2.8}
	For Problem \ref{prob:3.2.8}:
	\begin{enumerate}
		\item Let $g: \mathbb{R} \to \mathbb{R}$ be a strictly increasing monotonic function on the support of $X$. By the formal definition of strict increasing monotonicity, for any pair of real numbers $x, y$, the inequality relation is preserved under the application of the inverse operator:
		\begin{align*}
			g(x) \le y \iff x \le g^{-1}(y).
		\end{align*}
		Defining the transformed random variable $Y = g(X)$, we compute its cumulative distribution function at any point of the range $y \in \mathbb{R}$:
		\begin{align*}
			F_Y(y) = P(Y \le y) = P(g(X) \le y) = P(X \le g^{-1}(y)) = F_X\left( g^{-1}(y) \right).
		\end{align*}
		Since the composition of a right-step function ($F_X$) with an increasing bijective function ($g^{-1}$) produces a new step function with jump points exactly at $y_k = g(x_k)$, the preservation of the discrete structure and the theoretical identity is formally verified.

		\item Let $h: \mathbb{R} \to \mathbb{R}$ be a strictly decreasing monotonic function on the support of $X$, and define $Z = h(X)$.
		Upon applying a strictly decreasing transformation, the direction of the mathematical inequality is reversed in the preimage domain:
		\begin{align*}
			h(X) \le z \iff X \ge h^{-1}(z).
		\end{align*}
		Evaluating the cumulative probability for the variable $Z$:
		\begin{align*}
			F_Z(z) = P(Z \le z) = P(h(X) \le z) = P(X \ge h^{-1}(z)).
		\end{align*}
		To express the event $X \ge h^{-1}(z)$ in terms of the cumulative distribution function (which by definition accumulates strictly to the left as $P(X \le x)$), we use the probability of the complementary event and its left-hand limit:
		\begin{align*}
			P(X \ge h^{-1}(z)) &= 1 - P(X < h^{-1}(z)) \\
			&= 1 - \lim_{u \to (h^{-1}(z))^-} F_X(u) = 1 - F_X\left( \left[h^{-1}(z)\right]^- \right).
		\end{align*}
		In conclusion, for a decreasing transformation, the cumulative distribution function of the new discrete variable is algebraically governed by the exact analytical formula:
		\begin{align*}
			F_Z(z) = 1 - F_X\left( \left[h^{-1}(z)\right]^- \right).
		\end{align*}
	\end{enumerate}
\end{solucion}

\begin{solucion}
	\label{sol:3.2.9}
	For Problem \ref{prob:3.2.9}:
	\begin{enumerate}
		\item Let $X$ be a pure discrete random variable whose measurable support is $\mathcal{S}_X = \set{x_k}_{k \in \mathcal{K}}$, where the index $\mathcal{K}$ is at most countably infinite and each point concentrates probability $p_k = P(X = x_k) > 0$ with $\sum_k p_k = 1$.
		The cumulative distribution function at any point $x \in \mathbb{R}$ is written as the sum of the probability masses over the support points less than or equal to $x$:
		\begin{align*}
			F(x) = \sum_{x_k \le x} p_k.
		\end{align*}
		Now consider the unit Heaviside step function $H(t)$, which satisfies:
		\begin{align*}
			H(x - x_k) = \begin{cases}
				1 & \text{if } x - x_k \ge 0 \iff x \ge x_k, \\
				0 & \text{if } x - x_k < 0 \iff x < x_k.
			\end{cases}
		\end{align*}
		Multiplying each term of the infinite series by the indicator $H(x - x_k)$, every summand where $x < x_k$ is automatically annulled (since the step equals zero), while every summand where $x \ge x_k$ contributes its full mass $p_k \cdot 1 = p_k$:
		\begin{align*}
			\sum_{k \in \mathcal{K}} p_k\, H(x - x_k) = \sum_{k : x_k \le x} p_k \cdot 1 + \sum_{k : x_k > x} p_k \cdot 0 = \sum_{x_k \le x} p_k = F(x).
		\end{align*}
		Since, by the second Kolmogorov axiom, the series of probabilities is absolutely convergent ($\sum |p_k| = \sum p_k = 1 < \infty$), this series of Heaviside functions converges uniformly on all of $\mathbb{R}$, proving the discrete structural representation.

		\item In the sense of ordinary differential calculus (real analysis with respect to Lebesgue measure on $\mathbb{R}$), the complement of the countable support $\mathbb{R} \setminus \mathcal{S}_X$ is an open set of full measure, consisting of a countable union of disjoint intervals within which the step function $F(x)$ is locally and absolutely constant (horizontal plateaus).
		Consequently, the ordinary derivative on each plateau is identically zero:
		\begin{align*}
			F'(x) = 0 \quad \text{for all } x \notin \mathcal{S}_X.
		\end{align*}
		Since the cardinality of the discontinuity points $\mathcal{S}_X$ is countable (Lebesgue measure equal to zero, $\mu_L(\mathcal{S}_X) = 0$), it is concluded in real-variable analysis that $F'(x) = 0$ almost everywhere ($\text{a.e.}$).

		However, to preserve analytical coherence and rigor in engineering and theoretical physics without losing the total probability mass, the **theory of distributions (or Sobolev-Schwartz generalized functions)** is invoked. Within this topological framework, the distributional derivative of the unit Heaviside step function $H(t)$ is defined as the Dirac point-mass measure centered at the origin, $\delta(t)$, whose integral satisfies $\int_{-\infty}^\infty \delta(t) dt = 1$.
		By differentiating term by term the convergent series derived in the previous part in the sense of distributions, we obtain:
		\begin{align*}
			f(x) = \frac{d}{dx} F(x) = \frac{d}{dx} \left[ \sum_{k \in \mathcal{K}} p_k\, H(x - x_k) \right] = \sum_{k \in \mathcal{K}} p_k\, \frac{d}{dx} H(x - x_k) = \sum_{k \in \mathcal{K}} p_k\, \delta(x - x_k).
		\end{align*}
		This beautiful generalized identity reconciles continuous differential calculus with discrete probability, showing that the generalized density of a pure discrete random variable is a linear superposition of Dirac impulses weighted by their probability masses.
	\end{enumerate}
\end{solucion}

\begin{solucion}
	\label{sol:3.2.10}
	For Problem \ref{prob:3.2.10}:
	\begin{enumerate}
		\item Let $\mathcal{R} = (x_1, x_2] \times (y_1, y_2]$ be the half-open rectangle in the Cartesian plane with $x_1 < x_2$ and $y_1 < y_2$. The joint probability that the random vector $(X, Y)$ falls within this region is denoted by $P((X,Y) \in \mathcal{R})$.
		By the geometry of measurable sets in $\mathbb{R}^2$, we can express the infinite lower-left quadrant with vertex at $(x_2, y_2)$ as the disjoint union of four rectangular regions:
		\begin{align*}
			(-\infty, x_2] \times (-\infty, y_2] &= \mathcal{R} \cup \left[ (-\infty, x_1] \times (y_1, y_2] \right] \\
			&\quad \cup \left[ (x_1, x_2] \times (-\infty, y_1] \right] \cup \left[ (-\infty, x_1] \times (-\infty, y_1] \right].
		\end{align*}
		Evaluating the joint probability measure on the total outer quadrant (which by formal definition constitutes $F_{X,Y}(x_2, y_2)$):
		\begin{align*}
			F_{X,Y}(x_2, y_2) &= P((X,Y) \in \mathcal{R}) + P(X \le x_1, y_1 < Y \le y_2) \\
			&\quad + P(x_1 < X \le x_2, Y \le y_1) + P(X \le x_1, Y \le y_1).
		\end{align*}
		Decomposing the intermediate terms via differences of distribution functions:
		\begin{align*}
			P(X \le x_1, y_1 < Y \le y_2) &= P(X \le x_1, Y \le y_2) - P(X \le x_1, Y \le y_1) \\
			&= F_{X,Y}(x_1, y_2) - F_{X,Y}(x_1, y_1), \\[4pt]
			P(x_1 < X \le x_2, Y \le y_1) &= P(X \le x_2, Y \le y_1) - P(X \le x_1, Y \le y_1) \\
			&= F_{X,Y}(x_2, y_1) - F_{X,Y}(x_1, y_1).
		\end{align*}
		Substituting both algebraic identities into the original conservation equation and grouping terms:
		\begin{align*}
			F_{X,Y}(x_2, y_2) &= P((X,Y) \in \mathcal{R}) + \left[ F_{X,Y}(x_1, y_2) - F_{X,Y}(x_1, y_1) \right] \\
			&\quad + \left[ F_{X,Y}(x_2, y_1) - F_{X,Y}(x_1, y_1) \right] + F_{X,Y}(x_1, y_1) \\
			&= P((X,Y) \in \mathcal{R}) + F_{X,Y}(x_1, y_2) + F_{X,Y}(x_2, y_1) - F_{X,Y}(x_1, y_1).
		\end{align*}
		Finally solving for the probability mass contained inside the rectangle $\mathcal{R}$:
		\begin{align*}
			P(x_1 < X \le x_2, y_1 < Y \le y_2) = F_{X,Y}(x_2, y_2) - F_{X,Y}(x_1, y_2) - F_{X,Y}(x_2, y_1) + F_{X,Y}(x_1, y_1).
		\end{align*}
		This expression represents the mixed second difference or the $n$-th variation of the bidimensional cumulative operator.

		\item To prove the Fréchet-Hoeffding upper bound at any point $(x, y) \in \mathbb{R}^2$, note that the bivariate intersection event is trivially contained within each of the two marginal events separately:
		\begin{align*}
			\set{\omega \in \Omega : X(\omega) \le x \text{ and } Y(\omega) \le y} \subseteq \set{\omega \in \Omega : X(\omega) \le x}, \\
			\set{\omega \in \Omega : X(\omega) \le x \text{ and } Y(\omega) \le y} \subseteq \set{\omega \in \Omega : Y(\omega) \le y}.
		\end{align*}
		By the monotonicity of the probability measure, the probability of the subset cannot exceed the probability of either of the supersets that contain it:
		\begin{align*}
			F_{X,Y}(x, y) \le F_X(x) \quad \text{and} \quad F_{X,Y}(x, y) \le F_Y(y) \implies F_{X,Y}(x, y) \le \min\left( F_X(x), F_Y(y) \right).
		\end{align*}
		To prove the lower bound, we invoke Boole's Inequality for the union of complementary events: the sum of the probabilities of any two events relates to the probability of their intersection via $P(A \cap B) = P(A) + P(B) - P(A \cup B)$.
		Setting $A = \set{X \le x}$ and $B = \set{Y \le y}$:
		\begin{align*}
			F_{X,Y}(x, y) = P(A \cap B) = F_X(x) + F_Y(y) - P(X \le x \text{ or } Y \le y).
		\end{align*}
		Since the measure of the union of events in a probability space is axiomatically bounded by unity ($P(A \cup B) \le 1$), subtracting a number less than or equal to $1$ yields a lower bound when substituting the extreme value $1$:
		\begin{align*}
			F_{X,Y}(x, y) \ge F_X(x) + F_Y(y) - 1.
		\end{align*}
		Moreover, since by the first Kolmogorov axiom no cumulative probability can be strictly negative ($F_{X,Y}(x, y) \ge 0$), we combine both restrictions into a single maximum operator:
		\begin{align*}
			F_{X,Y}(x, y) \ge \max\left( 0, F_X(x) + F_Y(y) - 1 \right).
		\end{align*}
		In conclusion, the discrete joint cumulative distribution function (and that of any general type) remains invariably confined between the two universal Fréchet-Hoeffding extreme bounds:
		\begin{align*}
			\max\left( 0, F_X(x) + F_Y(y) - 1 \right) \le F_{X,Y}(x, y) \le \min\left( F_X(x), F_Y(y) \right).
		\end{align*}
	\end{enumerate}
\end{solucion}

% ==============================================================================
% SECTION 03.03: MATHEMATICAL EXPECTATION AND OPERATIONAL VARIANCE
% ==============================================================================

\subsection*{Problem Statements --- Section 03.03}

\subsubsection*{Fundamental Level}

\begin{problema}
	\label{prob:3.3.1}
	Let $X$ be a discrete random variable modeling the number of erroneous transactions during peak hour in a banking database, with support $\mathcal{S}_X = \set{0, 1, 2, 3, 4}$ and explicit probability mass function given by:
	\begin{table}[h]
		\centering
		\begin{tabular}{c|ccccc}
			\toprule
			$x$ & $0$ & $1$ & $2$ & $3$ & $4$ \\
			\midrule
			$f_X(x)$ & $0.35$ & $0.30$ & $0.20$ & $0.10$ & $0.05$ \\
			\bottomrule
		\end{tabular}
	\end{table}
	\begin{enumerate}
		\item Formally compute the mathematical expectation $\mu = \mathbb{E}[X]$ using the fundamental analytical definition.
		\item Compute the non-central second moment, or second-order raw moment, $\mathbb{E}[X^2]$.
		\item Evaluate the operational variance of $X$ using the classic König-Huygens identity ($\text{Var}(X) = \mathbb{E}[X^2] - \mu^2$) and verify its exact numerical equivalence by computing the sum of squared deviations from the mean via the central definition ($\sum (x-\mu)^2 f_X(x)$).
	\end{enumerate}
\end{problema}

\begin{problema}
	\label{prob:3.3.2}
	Consider a random experiment in which two perfectly balanced, independent six-sided dice are rolled simultaneously. Let us define the sum random variable $S = X_1 + X_2$, where $X_1$ and $X_2$ represent the faces obtained on the first and second die, respectively.
	\begin{enumerate}
		\item Compute the arithmetic mean of the support of $S$ (where $\mathcal{S}_S = \set{2, 3, \dots, 12}$) and rigorously explain why this simple arithmetic mean differs from, or coincides with, the numerical value of the true mathematical expectation $\mathbb{E}[S]$.
		\item Analytically compute the exact value of the expectation $\mathbb{E}[S]$ and the variance $\text{Var}(S)$ without tabulating the eleven values of the joint distribution of $S$, basing your development on the properties of linearity and independence.
	\end{enumerate}
\end{problema}

\begin{problema}
	\label{prob:3.3.3}
	Let $X$ be a discrete random variable whose support is the symmetric set $\mathcal{S}_X = \set{-2, -1, 0, 1, 2}$, with uniform probability mass function over said support ($f_X(x) = \frac{1}{5}$ for all $x \in \mathcal{S}_X$). Consider the non-bijective polynomial transformation given by:
	\begin{align*}
		g(X) = 2X^2 - 3X + 5.
	\end{align*}
	\begin{enumerate}
		\item Compute the expected value $\mathbb{E}[g(X)]$ by directly applying the Law of the Unconscious Statistician (LOTUS).
		\item Explicitly determine the induced probability distribution $Y = g(X)$ (its support $\mathcal{S}_Y$ and its mass function $f_Y(y)$) and formally verify the invariance of the result by evaluating $\sum y f_Y(y)$.
	\end{enumerate}
\end{problema}

\subsubsection*{Operational Level}

\begin{problema}
	\label{prob:3.3.4}
	A civil engineering and infrastructure corporation evaluates the bidding of a mining megaproject whose net profitability in millions of dollars is a discrete random variable $R$ with the following distribution conditioned on global economic scenarios:
	\begin{table}[h]
		\centering
		\begin{tabular}{l|ccc}
			\toprule
			\textbf{Economic Scenario} & \textbf{Probability ($p_k$)} & \textbf{Net Profitability ($r_k$)} \\
			\midrule
			Accelerated expansion & $0.25$ & $+120$ M USD \\
			Moderate stability & $0.55$ & $+40$ M USD \\
			Severe recession & $0.20$ & $-50$ M USD \\
			\bottomrule
		\end{tabular}
	\end{table}
	\begin{enumerate}
		\item Determine the expected net present value of the project $\mathbb{E}[R]$ and the standard deviation of the return $\sigma_R$, physically interpreting the latter as a measure of capital volatility or risk.
		\item If management requires a minimum risk-adjusted return defined by the financial safety metric $U = \mathbb{E}[R] - 1.25\sigma_R$, should the corporation accept the bid if the minimum approval threshold is $U \ge 0$?
	\end{enumerate}
\end{problema}

\begin{problema}
	\label{prob:3.3.5}
	A cloud data processing center manages a massive cluster of servers. Let $K$ be the number of servers that experience a critical outage or failure during a 24-hour operating window. Historical statistical studies establish that $\mathbb{E}[K] = 3.2$ servers down and that the operational variance is $\text{Var}(K) = 1.8$.
	The total operating cost of repair, SLA contract penalty, and load redistribution (in thousands of dollars) follows a strictly quadratic and strictly convex polynomial structure in the number of failures:
	\begin{align*}
		C(K) = 15 + 12K + 4K^2.
	\end{align*}
	Compute exactly the expected operating cost $\mathbb{E}[C(K)]$ that management must budget per day of operation without needing to know the complete probability distribution of $K$.
\end{problema}

\begin{problema}
	\label{prob:3.3.6}
	In a global technical assessment for industrial cybersecurity certification, the raw score $X$ of the evaluated engineers has mean $\mu_X = 68.5$ points and standard deviation $\sigma_X = 12.4$ points. To standardize results among different international committees, the standardized random variable (or $Z$-score) is defined as:
	\begin{align*}
		Z = \frac{X - \mu_X}{\sigma_X}.
	\end{align*}
	Formally prove, with detailed algebraic steps, that for any discrete random variable $X$ with $\sigma_X > 0$, the standardized variable $Z$ rigorously satisfies that its mathematical expectation is zero ($\mathbb{E}[Z] = 0$) and its variance is unitary ($\text{Var}(Z) = 1$).
\end{problema}

\subsubsection*{Analytical Level}

\begin{problema}
	\label{prob:3.3.7}
	Let $X$ be a discrete random variable with finite mean $\mu$ and finite variance $\sigma^2$. Consider the mean-squared-error function with respect to an arbitrary scalar $a \in \mathbb{R}$, defined as the second-order moment centered at $a$:
	\begin{align*}
		M(a) = \mathbb{E}\left[ (X - a)^2 \right].
	\end{align*}
	\begin{enumerate}
		\item Rigorously prove, by algebraic expansion and use of the linearity of the expectation operator, the fundamental moment-decomposition identity:
		\begin{align*}
			M(a) = \text{Var}(X) + (\mu - a)^2.
		\end{align*}
		\item From the proven identity, exactly derive the unique value of $a$ that minimizes the mean squared error of the probabilistic approximation, proving the classical theorem of the probabilistic center of gravity:
		\begin{align*}
			\min_{a \in \mathbb{R}} \mathbb{E}\left[ (X - a)^2 \right] = \text{Var}(X).
		\end{align*}
	\end{enumerate}
\end{problema}

\begin{problema}
	\label{prob:3.3.8}
	Consider two discrete random variables $X$ and $Y$ defined on the same probability space $\Omega$.
	\begin{enumerate}
		\item Rigorously prove, from the joint summation operator and the axiomatic properties of the marginal laws, that for any pair of scalars $c, d \in \mathbb{R}$ the universal linearity holds:
		\begin{align*}
			\mathbb{E}[cX + dY] = c\mathbb{E}[X] + d\mathbb{E}[Y],
		\end{align*}
		regardless of whether the variables $X$ and $Y$ are independent or dependent.
		\item Analytically prove that if $X$ and $Y$ are strictly and functionally independent ($\mathbb{E}[XY] = \mathbb{E}[X]\mathbb{E}[Y]$), then the variance operator satisfies the second-degree additive homogeneity property:
		\begin{align*}
			\text{Var}(cX \pm dY) = c^2\text{Var}(X) + d^2\text{Var}(Y).
		\end{align*}
	\end{enumerate}
\end{problema}

\subsubsection*{Challenging Level}

\begin{problema}
	\label{prob:3.3.9}
	Consider the famous random game known as the **St. Petersburg Paradox**: in an ideal casino, a fair coin is tossed consecutively until the first heads appears on the $k$-th toss. At that exact instant the game ends and the casino pays the player a monetary prize of $X = 2^k$ monetary units.
	\begin{enumerate}
		\item Formally prove that the mathematical expectation of the prize $\mathbb{E}[X]$ diverges to infinity ($\mathbb{E}[X] = \infty$), evidencing the economic paradox between an infinite expected value and a rational player's unwillingness to pay a substantial sum to participate.
		\item **Regularization via the casino's bankroll limit:** If, by legal regulation, the casino imposes a maximum capital reserve cap equivalent to $N$ tosses, so that if heads does not appear within the first $N$ attempts the game stops and the truncated maximum prize $M = 2^N$ is paid, determine a closed-form expression for the truncated expected value $\mathbb{E}[X_N]$. Compute the exact numerical value that a player should pay if the house's bankroll limit is $N = 25$ (nearly 33.5 million coins).
	\end{enumerate}
\end{problema}

\begin{problema}
	\label{prob:3.3.10}
	Let $X$ be any discrete random variable whose support $\mathcal{S}_X$ is strictly bounded within a closed interval of the real line:
	\begin{align*}
		P(a \le X \le b) = 1, \quad \text{where } -\infty < a < b < \infty.
	\end{align*}
	Formally prove the **Popoviciu Supremum Bound Theorem for Discrete Variance**, which establishes that the maximum variance attainable by any variable with support bounded in $[a, b]$ is strictly bounded by:
	\begin{align*}
		\text{Var}(X) \le \frac{(b - a)^2}{4}.
	\end{align*}
	Additionally specify which discrete distribution attains this absolute upper bound with exact equality.
\end{problema}

% ==============================================================================
% HINTS FOR SECTION 03.03
% ==============================================================================

\subsection*{Hints for the Problems --- Section 03.03}

\begin{sugerencia}
	\label{sug:3.3.1}
	For Problem \ref{prob:3.3.1}, organize the calculations in an auxiliary table adding columns for $x \cdot f_X(x)$, $x^2 \cdot f_X(x)$, and $(x-\mu)^2 \cdot f_X(x)$. Recall that in floating-point or fractional arithmetic the König-Huygens identity is an invaluable algebraic shortcut.
\end{sugerencia}

\begin{sugerencia}
	\label{sug:3.3.2}
	For Problem \ref{prob:3.3.2}, it is not necessary to tabulate the 11 possible values of the sum. Directly apply the linearity operator $\mathbb{E}[S] = \mathbb{E}[X_1] + \mathbb{E}[X_2]$ and the additivity of variance for independent variables $\text{Var}(S) = \text{Var}(X_1) + \text{Var}(X_2)$, recalling that for a single die $\mathbb{E}[X_i] = 3.5$ and $\text{Var}(X_i) = \frac{35}{12}$.
\end{sugerencia}

\begin{sugerencia}
	\label{sug:3.3.3}
	For Problem \ref{prob:3.3.3}, notice that the function $g(X) = 2X^2 - 3X + 5$ is not one-to-one (for example, $g(-1) = 2(1)+3+5=10$ and $g(1) = 2-3+5=4$, but $g(-2) = 19$ and $g(2) = 7$). To verify LOTUS, group the probabilities of those $x$ that produce the same image at $y = g(x)$.
\end{sugerencia}

\begin{sugerencia}
	\label{sug:3.3.4}
	For Problem \ref{prob:3.3.4}, first compute the expectation as the probability-weighted average $\mathbb{E}[R] = \sum r_k p_k$. Then find the variance by summing the weighted squared deviations $\sum (r_k - \mathbb{E}[R])^2 p_k$ and take its positive square root to obtain $\sigma_R$.
\end{sugerencia}

\begin{sugerencia}
	\label{sug:3.3.5}
	For Problem \ref{prob:3.3.5}, expand the linear operator over the cost: $\mathbb{E}[C(K)] = 15 + 12\mathbb{E}[K] + 4\mathbb{E}[K^2]$. To obtain the non-central second moment $\mathbb{E}[K^2]$, solve from the classic definition of variance: $\mathbb{E}[K^2] = \text{Var}(K) + (\mathbb{E}[K])^2$.
\end{sugerencia}

\begin{sugerencia}
	\label{sug:3.3.6}
	For Problem \ref{prob:3.3.6}, factor out the constants $1/\sigma_X$ and $\mu_X/\sigma_X$ from the linear expectation operator. For the variance, use the fundamental invariance property under translation and quadratic scaling: $\text{Var}(aX + b) = a^2\text{Var}(X)$.
\end{sugerencia}

\begin{sugerencia}
	\label{sug:3.3.7}
	For Problem \ref{prob:3.3.7}, add and subtract the true mean $\mu$ inside the squared parenthesis: $(X - a)^2 = [(X - \mu) + (\mu - a)]^2$. When expanding the squared binomial and applying the expectation term by term, observe what happens to the cross term $2(\mu - a)\mathbb{E}[X - \mu]$.
\end{sugerencia}

\begin{sugerencia}
	\label{sug:3.3.8}
	For Problem \ref{prob:3.3.8}, in the second part algebraically expand the definition of the central moment: $\text{Var}(cX + dY) = \mathbb{E}[((cX + dY) - (c\mu_X + d\mu_Y))^2]$. Group by variables $[c(X-\mu_X) + d(Y-\mu_Y)]^2$ and recall that independence implies the covariance is strictly zero ($\mathbb{E}[(X-\mu_X)(Y-\mu_Y)] = 0$).
\end{sugerencia}

\begin{sugerencia}
	\label{sug:3.3.9}
	For Problem \ref{prob:3.3.9}, in the first part set up the infinite series of the expectation by summing the terms of the product of the payoff by the probability: $\sum_{k=1}^\infty (2^k)\left(\frac{1}{2^k}\right)$. In the second part, sum the first $N$ terms of the progression and analyze the exact algebraic behavior of the truncated expectation as a function of the parameter $N$.
\end{sugerencia}

\begin{sugerencia}
	\label{sug:3.3.10}
	For Problem \ref{prob:3.3.10}, start from the non-negative algebraic identity $(X - a)(b - X) \ge 0$, valid with probability $1$ given that the support is bounded in $[a, b]$. Take expectation on both sides and expand the product, relating $\mathbb{E}[X^2]$ to $\mathbb{E}[X]$ and completing the square with respect to the mean $\mu$.
\end{sugerencia}

% ==============================================================================
% SOLUTIONS FOR SECTION 03.03
% ==============================================================================

\subsection*{Solutions to the Problems --- Section 03.03}

\begin{solucion}
	\label{sol:3.3.1}
	For Problem \ref{prob:3.3.1}:
	\begin{enumerate}
		\item The analytical computation of the mathematical expectation or mean $\mu = \mathbb{E}[X]$ is performed by weighting each element of the support by its respective probability mass:
		\begin{align*}
			\mu &= \sum_{x \in \mathcal{S}_X} x \cdot f_X(x) \\
			&= (0)(0.35) + (1)(0.30) + (2)(0.20) + (3)(0.10) + (4)(0.05) \\
			&= 0.00 + 0.30 + 0.40 + 0.30 + 0.20 \\
			&= 1.20 \text{ erroneous transactions}.
		\end{align*}
		Physical-statistical meaning: In a representative peak hour, the long-run weighted average number of database errors is exactly $1.20$.

		\item For the raw second-order moment $\mathbb{E}[X^2]$, we weight the exact squares of the abscissas of the support:
		\begin{align*}
			\mathbb{E}[X^2] &= \sum_{x \in \mathcal{S}_X} x^2 \cdot f_X(x) \\
			&= (0^2)(0.35) + (1^2)(0.30) + (2^2)(0.20) + (3^2)(0.10) + (4^2)(0.05) \\
			&= (0)(0.35) + (1)(0.30) + (4)(0.20) + (9)(0.10) + (16)(0.05) \\
			&= 0.00 + 0.30 + 0.80 + 0.90 + 0.80 \\
			&= 2.80.
		\end{align*}

		\item **Variance calculation via the König-Huygens formula:**
		We apply the fundamental identity connecting the raw moments to the central variance:
		\begin{align*}
			\text{Var}(X) &= \mathbb{E}[X^2] - \mu^2 \\
			&= 2.80 - (1.20)^2 \\
			&= 2.80 - 1.44 = 1.36 \text{ (erroneous transactions)}^2.
		\end{align*}
		**Cross-check using the classic central-variance definition:**
		We evaluate the explicit sum of the squared deviations weighted by their probability:
		\begin{align*}
			\text{Var}(X) &= \sum_{x \in \mathcal{S}_X} (x - \mu)^2 \cdot f_X(x) \\
			&= (0 - 1.20)^2(0.35) + (1 - 1.20)^2(0.30) + (2 - 1.20)^2(0.20) \\
			&\quad + (3 - 1.20)^2(0.10) + (4 - 1.20)^2(0.05) \\
			&= (1.44)(0.35) + (0.04)(0.30) + (0.64)(0.20) + (3.24)(0.10) + (7.84)(0.05) \\
			&= 0.504 + 0.012 + 0.128 + 0.324 + 0.392 \\
			&= 1.36 \text{ (erroneous transactions)}^2.
		\end{align*}
		The numerical equivalence between both methods ($1.36$) validates the algebraic exactness of the computation.
	\end{enumerate}
\end{solucion}

\begin{solucion}
	\label{sol:3.3.2}
	For Problem \ref{prob:3.3.2}:
	\begin{enumerate}
		\item The support of the sum of two dice is $\mathcal{S}_S = \set{2, 3, 4, 5, 6, 7, 8, 9, 10, 11, 12}$, made up of $N = 11$ discrete elements.
		The simple, unweighted arithmetic mean of the support set is calculated as:
		\begin{align*}
			\bar{S}_{\text{support}} = \frac{1}{N} \sum_{k=2}^{12} k = \frac{2 + 3 + 4 + 5 + 6 + 7 + 8 + 9 + 10 + 11 + 12}{11} = \frac{77}{11} = 7.
		\end{align*}
		In this extraordinarily particular and symmetric case, the numerical value coincides with the true center of gravity or mathematical expectation ($\mathbb{E}[S] = 7$). However, it is vital to rigorously clarify that **both concepts fundamentally differ by definition**: the arithmetic mean of the support assumes uniform weighting ($1/11$) for all values of the range, whereas the true mathematical expectation weights each sum $k$ by its true, non-uniform triangular probability ($f_S(k) = \frac{6 - |7-k|}{36}$). The coincidence is due exclusively to the fact that the probability distribution $f_S(k)$ is symmetric about its geometric center $k=7$.

		\item For the analytical calculation without combinatorial tabulation, we invoke the universal linearity of mathematical expectation and the functional independence of the dice:
		\begin{align*}
			\mathbb{E}[S] = \mathbb{E}[X_1 + X_2] = \mathbb{E}[X_1] + \mathbb{E}[X_2].
		\end{align*}
		Since each individual die $X_i$ is an equiprobable random variable with support $\set{1, 2, 3, 4, 5, 6}$, its elementary expectation is:
		\begin{align*}
			\mathbb{E}[X_i] = \sum_{j=1}^6 j \cdot \frac{1}{6} = \frac{1+2+3+4+5+6}{6} = \frac{21}{6} = 3.5.
		\end{align*}
		Therefore:
		\begin{align*}
			\mathbb{E}[S] = 3.5 + 3.5 = 7.
		\end{align*}
		For the variance, since the tosses are strictly independent events on a product space, the covariance is zero ($\text{Cov}(X_1, X_2) = 0$) and the additive variance simplifies to:
		\begin{align*}
			\text{Var}(S) = \text{Var}(X_1 + X_2) = \text{Var}(X_1) + \text{Var}(X_2).
		\end{align*}
		We compute the elementary variance of a single die via $E[X_i^2] - (\mathbb{E}[X_i])^2$:
		\begin{align*}
			\mathbb{E}[X_i^2] = \sum_{j=1}^6 j^2 \cdot \frac{1}{6} = \frac{1 + 4 + 9 + 16 + 25 + 36}{6} = \frac{91}{6}.
		\end{align*}
		Then, the variance of each face is:
		\begin{align*}
			\text{Var}(X_i) = \frac{91}{6} - \left(\frac{7}{2}\right)^2 = \frac{91}{6} - \frac{49}{4} = \frac{182 - 147}{12} = \frac{35}{12}.
		\end{align*}
		Summing both marginal variances:
		\begin{align*}
			\text{Var}(S) = \frac{35}{12} + \frac{35}{12} = \frac{70}{12} = \frac{35}{6} \approx 5.8333.
		\end{align*}
	\end{enumerate}
\end{solucion}

\begin{solucion}
	\label{sol:3.3.3}
	For Problem \ref{prob:3.3.3}:
	\begin{enumerate}
		\item The support of $X$ is $\mathcal{S}_X = \set{-2, -1, 0, 1, 2}$ and $f_X(x) = \frac{1}{5}$ constant for the 5 points.
		We directly apply the Law of the Unconscious Statistician (LOTUS), evaluating the sum of the images of $g(x)$ multiplied by the original masses:
		\begin{align*}
			\mathbb{E}[g(X)] &= \sum_{x \in \mathcal{S}_X} g(x) \cdot f_X(x) = \sum_{x \in \set{-2, -1, 0, 1, 2}} (2x^2 - 3x + 5) \cdot \frac{1}{5} \\
			&= \frac{1}{5} \left[ g(-2) + g(-1) + g(0) + g(1) + g(2) \right].
		\end{align*}
		We numerically evaluate each polynomial image:
		\begin{align*}
			g(-2) &= 2(-2)^2 - 3(-2) + 5 = 2(4) + 6 + 5 = 19, \\
			g(-1) &= 2(-1)^2 - 3(-1) + 5 = 2(1) + 3 + 5 = 10, \\
			g(0) &= 2(0)^2 - 3(0) + 5 = 0 + 0 + 5 = 5, \\
			g(1) &= 2(1)^2 - 3(1) + 5 = 2(1) - 3 + 5 = 4, \\
			g(2) &= 2(2)^2 - 3(2) + 5 = 2(4) - 6 + 5 = 7.
		\end{align*}
		Substituting the evaluations into the sum:
		\begin{align*}
			\mathbb{E}[g(X)] = \frac{1}{5} (19 + 10 + 5 + 4 + 7) = \frac{45}{5} = 9.
		\end{align*}

		\item To construct the induced mass distribution $Y = g(X)$, we identify the values of the range or image taken by the transformation:
		$\mathcal{S}_Y = \set{4, 5, 7, 10, 19}$.
		Since the five images obtained ($19, 10, 5, 4, 7$) are all completely distinct from each other (there was no collision due to symmetry in this particular parabolic quadratic because of the linear term $-3X$), each point of the induced support inherits exactly one unique preimage of mass $1/5$:
		\begin{table}[h]
			\centering
			\begin{tabular}{c|ccccc}
				\toprule
				$y$ & $4$ & $5$ & $7$ & $10$ & $19$ \\
				\midrule
				$f_Y(y)$ & $0.20$ & $0.20$ & $0.20$ & $0.20$ & $0.20$ \\
				\bottomrule
			\end{tabular}
		\end{table}

		We verify the expectation directly in the new induced probability space:
		\begin{align*}
			\mathbb{E}[Y] = \sum_{y \in \mathcal{S}_Y} y \cdot f_Y(y) = (4+5+7+10+19) \cdot \frac{1}{5} = \frac{45}{5} = 9.
		\end{align*}
		Both procedures conclude identically at the exact expected value of $9$, illustrating the great computational savings provided by the LOTUS theorem by avoiding the formal construction of $f_Y(y)$.
	\end{enumerate}
\end{solucion}

\begin{solucion}
	\label{sol:3.3.4}
	For Problem \ref{prob:3.3.4}:
	\begin{enumerate}
		\item We first compute the expected value of the net profitability $\mathbb{E}[R]$ by summing the products of each conditional return by its macroeconomic probability:
		\begin{align*}
			\mathbb{E}[R] &= \sum_{k=1}^3 r_k \cdot p_k \\
			&= (120)(0.25) + (40)(0.55) + (-50)(0.20) \\
			&= 30.0 + 22.0 - 10.0 = +42.0 \text{ Million USD}.
		\end{align*}
		To compute the volatility or standard deviation $\sigma_R$, we first evaluate the raw second moment of the profitability $\mathbb{E}[R^2]$:
		\begin{align*}
			\mathbb{E}[R^2] &= \sum_{k=1}^3 r_k^2 \cdot p_k \\
			&= (120^2)(0.25) + (40^2)(0.55) + ((-50)^2)(0.20) \\
			&= (14400)(0.25) + (1600)(0.55) + (2500)(0.20) \\
			&= 3600 + 880 + 500 = 4980 \text{ (M USD)}^2.
		\end{align*}
		Applying the König-Huygens computational identity for the variance:
		\begin{align*}
			\text{Var}(R) &= \mathbb{E}[R^2] - (\mathbb{E}[R])^2 \\
			&= 4980 - (42.0)^2 = 4980 - 1764 = 3216 \text{ (M USD)}^2.
		\end{align*}
		Extracting the square root, we obtain the standard deviation of the return:
		\begin{align*}
			\sigma_R = \sqrt{3216} \approx 56.7098 \text{ Million USD}.
		\end{align*}
		Financial engineering interpretation: Although the positive expected return of $+42.0$ M USD is highly attractive, the enormous standard deviation of $\approx 56.71$ M USD reflects extreme capital volatility, alerting the board to the real danger of suffering catastrophic losses (up to $-50$ M USD) in the event of a recessionary shock in the mining market.

		\item We evaluate the corporate decision threshold or risk-adjusted utility $U$:
		\begin{align*}
			U &= \mathbb{E}[R] - 1.25\sigma_R \\
			&= 42.0 - 1.25(56.7098) = 42.0 - 70.88725 = -28.88725 \text{ M USD}.
		\end{align*}
		Since the risk-adjusted utility metric yields a strictly negative value ($U \approx -28.89 < 0$), **the corporation must formally reject the bid for the megaproject**. The expected profitability surplus fails to compensate or protect the company against the enormous financial risk involved.
	\end{enumerate}
\end{solucion}

\begin{solucion}
	\label{sol:3.3.5}
	For Problem \ref{prob:3.3.5}:
	The total maintenance and repair cost is a second-degree polynomial function of the number of downed servers $K$:
	\begin{align*}
		C(K) = 15 + 12K + 4K^2.
	\end{align*}
	By the linear-operator property of mathematical expectation (Theorem \ref{thm:2.9.1}), we apply the operator to each term of the linear combination:
	\begin{align*}
		\mathbb{E}[C(K)] &= \mathbb{E}[15 + 12K + 4K^2] \\
		&= 15 + 12\mathbb{E}[K] + 4\mathbb{E}[K^2].
	\end{align*}
	The statement directly provides the mean number of failures ($\mathbb{E}[K] = 3.2$) and the operational variance ($\text{Var}(K) = 1.8$). To evaluate the quadratic term $\mathbb{E}[K^2]$, we solve from the classic variance formula:
	\begin{align*}
		\text{Var}(K) = \mathbb{E}[K^2] - (\mathbb{E}[K])^2 \implies \mathbb{E}[K^2] = \text{Var}(K) + (\mathbb{E}[K])^2.
	\end{align*}
	Substituting the cluster's numerical values:
	\begin{align*}
		\mathbb{E}[K^2] = 1.8 + (3.2)^2 = 1.8 + 10.24 = 12.04.
	\end{align*}
	Finally, we substitute the moments $\mathbb{E}[K]$ and $\mathbb{E}[K^2]$ into the linear expansion of the expected cost:
	\begin{align*}
		\mathbb{E}[C(K)] &= 15 + 12(3.2) + 4(12.04) \\
		&= 15 + 38.4 + 48.16 \\
		&= 101.56 \text{ thousand dollars}.
	\end{align*}
	In conclusion, the finance department must reserve an exact average budget of **\$101,560 USD per day** to cover contingency management at the data center without needing to know or assume a parametric distribution for $K$.
\end{solucion}

\begin{solucion}
	\label{sol:3.3.6}
	For Problem \ref{prob:3.3.6}:
	Consider the affine standardization transformation $Z = \frac{X - \mu_X}{\sigma_X}$. We can separate this fraction into the sum of a constant slope $a = \frac{1}{\sigma_X}$ and a constant intercept $b = -\frac{\mu_X}{\sigma_X}$:
	\begin{align*}
		Z = aX + b, \quad \text{where } a = \frac{1}{\sigma_X} \text{ and } b = -\frac{\mu_X}{\sigma_X}.
	\end{align*}

	**Proof that the standardized expectation is zero ($\mathbb{E}[Z] = 0$):**
	We apply the linearity of the expectation operator to the affine combination:
	\begin{align*}
		\mathbb{E}[Z] &= \mathbb{E}[aX + b] = a\mathbb{E}[X] + b \\
		&= \left(\frac{1}{\sigma_X}\right) \mathbb{E}[X] - \frac{\mu_X}{\sigma_X}.
	\end{align*}
	By the fundamental definition of the mean of the original support, we know that $\mathbb{E}[X] = \mu_X$. Substituting this identity produces an exact cancellation of opposite terms:
	\begin{align*}
		\mathbb{E}[Z] = \frac{\mu_X}{\sigma_X} - \frac{\mu_X}{\sigma_X} = 0.
	\end{align*}

	**Proof that the standardized variance is unitary ($\text{Var}(Z) = 1$):**
	We invoke the two laws of variance transformation under linear changes (Theorem \ref{thm:2.9.3} of the master textbook):
	1. Variance is strictly invariant under translations or constant shifts ($\text{Var}(Y + c) = \text{Var}(Y)$ for all $c \in \mathbb{R}$). Hence, the constant term $b$ vanishes.
	2. When multiplying by a scalar $a$, the variance scales quadratically by that factor ($\text{Var}(aY) = a^2\text{Var}(Y)$).

	Applying both algebraic laws to $Z$:
	\begin{align*}
		\text{Var}(Z) &= \text{Var}(aX + b) = \text{Var}(aX) \\
		&= a^2 \text{Var}(X) = \left(\frac{1}{\sigma_X}\right)^2 \text{Var}(X) \\
		&= \frac{\text{Var}(X)}{\sigma_X^2}.
	\end{align*}
	Since by definition the standard deviation is the positive square root of the variance ($\sigma_X^2 \equiv \text{Var}(X)$), the fraction results in an exact identity with identical numerator and denominator:
	\begin{align*}
		\text{Var}(Z) = \frac{\sigma_X^2}{\sigma_X^2} = 1.
	\end{align*}
	This is rigorously and generally proven for any discrete or continuous random variable with non-degenerate dispersion.
\end{solucion}

\begin{solucion}
	\label{sol:3.3.7}
	For Problem \ref{prob:3.3.7}:
	\begin{enumerate}
		\item We want to expand and simplify the mean-squared-error function $M(a) = \mathbb{E}[(X - a)^2]$.
		To directly connect this centered moment with the central invariants of the system, we introduce an artificial zero into the difference by subtracting and adding the true population mean $\mu \equiv \mathbb{E}[X]$:
		\begin{align*}
			X - a = (X - \mu) + (\mu - a).
		\end{align*}
		We square both sides of the equality, grouping the probabilistic term $(X-\mu)$ on one hand and the constant difference $(\mu-a)$ on the other:
		\begin{align*}
			(X - a)^2 &= \left[ (X - \mu) + (\mu - a) \right]^2 \\
			&= (X - \mu)^2 + 2(X - \mu)(\mu - a) + (\mu - a)^2.
		\end{align*}
		Now, we take mathematical expectation on both extremes, invoking the universal linearity of the operator:
		\begin{align*}
			M(a) = \mathbb{E}\left[ (X - a)^2 \right] &= \mathbb{E}\left[ (X - \mu)^2 \right] + \mathbb{E}\left[ 2(X - \mu)(\mu - a) \right] + \mathbb{E}\left[ (\mu - a)^2 \right].
		\end{align*}
		Let us carefully analyze the three resulting terms:
		- The first term is, by the canonical definition of central dispersion, exactly the variance: $\mathbb{E}[(X - \mu)^2] = \text{Var}(X)$.
		- In the second term, the quantities $2$ and $(\mu - a)$ are exact numerical constants that come out of the linear expectation operator:
		  \begin{align*}
		  	\mathbb{E}\left[ 2(X - \mu)(\mu - a) \right] = 2(\mu - a) \mathbb{E}[X - \mu].
		  \end{align*}
		  However, the expectation of a random variable centered at its own mean is always identically zero ($\mathbb{E}[X - \mu] = \mathbb{E}[X] - \mu = \mu - \mu = 0$). Consequently, the entire cross covariance term vanishes: $2(\mu - a)(0) = 0$.
		- The third term is the average of a pure invariant constant, so it remains identical: $\mathbb{E}[(\mu - a)^2] = (\mu - a)^2$.

		Bringing together the nonzero algebraic pieces, we obtain the exact identity of the squared error:
		\begin{align*}
			M(a) = \text{Var}(X) + (\mu - a)^2.
		\end{align*}

		\item To find the scalar $a \in \mathbb{R}$ that minimizes $M(a)$, we observe the analytical structure of our decomposition:
		$\text{Var}(X)$ is a structural constant that depends solely on the intrinsic dispersion of $X$ (hence, it is immutable with respect to $a$).
		The second term $(\mu - a)^2$ is the square of a difference of real numbers, which guarantees, by the order axioms in $\mathbb{R}$, that it is a non-negative quantity:
		\begin{align*}
			(\mu - a)^2 \ge 0 \quad \text{for all } a \in \mathbb{R}.
		\end{align*}
		As a direct consequence, adding the variance to both sides, we obtain the universal bounded inequality:
		\begin{align*}
			M(a) = \text{Var}(X) + (\mu - a)^2 \ge \text{Var}(X) + 0 = \text{Var}(X).
		\end{align*}
		The absolute lower bound $\text{Var}(X)$ is achieved with exact equality if, and only if, the non-negative quadratic term vanishes completely:
		\begin{align*}
			(\mu - a)^2 = 0 \iff a = \mu.
		\end{align*}
		Consequently, the Center-of-Gravity Theorem is rigorously proven:
		\begin{align*}
			\min_{a \in \mathbb{R}} \mathbb{E}\left[ (X - a)^2 \right] = \text{Var}(X), \quad \text{attained only when } a = \mu = \mathbb{E}[X].
		\end{align*}
	\end{enumerate}
\end{solucion}

\begin{solucion}
	\label{sol:3.3.8}
	For Problem \ref{prob:3.3.8}:
	\begin{enumerate}
		\item Let $X$ and $Y$ be two discrete random variables with a joint probability mass function given by $f_{X,Y}(x, y) = P(X = x, Y = y)$ for every pair $(x, y)$ in the Cartesian product of the supports $\mathcal{S}_X \times \mathcal{S}_Y$.
		By the fundamental analytical definition in the multivariate space, the expected value of any affine linear combination $cX + dY$ is computed by summing over the entire two-dimensional joint support:
		\begin{align*}
			\mathbb{E}[cX + dY] = \sum_{x \in \mathcal{S}_X} \sum_{y \in \mathcal{S}_Y} (cx + dy) \cdot f_{X,Y}(x, y).
		\end{align*}
		We apply the algebraic distributive property of the product with respect to the sum in each term:
		\begin{align*}
			\mathbb{E}[cX + dY] = \sum_{x \in \mathcal{S}_X} \sum_{y \in \mathcal{S}_Y} \left[ cx \cdot f_{X,Y}(x, y) + dy \cdot f_{X,Y}(x, y) \right].
		\end{align*}
		By the elementary properties of absolute convergence in summable discrete series, we can separate the double sum into the addition of two independent two-dimensional series:
		\begin{align*}
			\mathbb{E}[cX + dY] = \sum_{x \in \mathcal{S}_X} \sum_{y \in \mathcal{S}_Y} cx \cdot f_{X,Y}(x, y) + \sum_{x \in \mathcal{S}_X} \sum_{y \in \mathcal{S}_Y} dy \cdot f_{X,Y}(x, y).
		\end{align*}
		In the first double sum, the factor $cx$ is algebraically independent of the summation index $y$, allowing it to be factored out of the inner sum. Similarly, in the second sum, the factor $dy$ is independent of the index $x$, so we swap the order of summation by the Fubini-Tonelli Theorem and factor out $dy$:
		\begin{align*}
			\mathbb{E}[cX + dY] = c \sum_{x \in \mathcal{S}_X} x \left[ \sum_{y \in \mathcal{S}_Y} f_{X,Y}(x, y) \right] + d \sum_{y \in \mathcal{S}_Y} y \left[ \sum_{x \in \mathcal{S}_X} f_{X,Y}(x, y) \right].
		\end{align*}
		By the marginal-consistency axiom of multivariate probability, summing or marginalizing the joint distribution with respect to one of the variables yields exactly the marginal mass function of the complementary variable:
		\begin{align*}
			\sum_{y \in \mathcal{S}_Y} f_{X,Y}(x, y) = f_X(x) \quad \text{and} \quad \sum_{x \in \mathcal{S}_X} f_{X,Y}(x, y) = f_Y(y).
		\end{align*}
		Substituting these marginal identities into our sums:
		\begin{align*}
			\mathbb{E}[cX + dY] &= c \sum_{x \in \mathcal{S}_X} x \cdot f_X(x) + d \sum_{y \in \mathcal{S}_Y} y \cdot f_Y(y) \\
			&= c\mathbb{E}[X] + d\mathbb{E}[Y].
		\end{align*}
		It is thereby shown that linearity is an absolute mathematical truth that holds universally, regardless entirely of the independence or correlation structure of $X$ and $Y$.

		\item Consider now the variance operator applied to the linear combination $W = cX \pm dY$.
		By the formal definition of central variance:
		\begin{align*}
			\text{Var}(W) = \mathbb{E}\left[ (W - \mathbb{E}[W])^2 \right].
		\end{align*}
		Using the linearity proved in part 1, the mean of our affine combination is $\mu_W = c\mu_X \pm d\mu_Y$ (where $\mu_X = \mathbb{E}[X]$ and $\mu_Y = \mathbb{E}[Y]$). Substituting into the central difference and grouping by the original variables:
		\begin{align*}
			W - \mu_W &= (cX \pm dY) - (c\mu_X \pm d\mu_Y) \\
			&= c(X - \mu_X) \pm d(Y - \mu_Y).
		\end{align*}
		We square the resulting binomial:
		\begin{align*}
			(W - \mu_W)^2 &= \left[ c(X - \mu_X) \pm d(Y - \mu_Y) \right]^2 \\
			&= c^2(X - \mu_X)^2 \pm 2cd(X - \mu_X)(Y - \mu_Y) + d^2(Y - \mu_Y)^2.
		\end{align*}
		We apply the linear expectation operator to each of the three expanded terms:
		\begin{align*}
			\text{Var}(cX \pm dY) = c^2 \mathbb{E}\left[ (X - \mu_X)^2 \right] \pm 2cd \mathbb{E}\left[ (X - \mu_X)(Y - \mu_Y) \right] + d^2 \mathbb{E}\left[ (Y - \mu_Y)^2 \right].
		\end{align*}
		We immediately recognize the definitions of the first and third summands as the marginal variances: $\mathbb{E}[(X-\mu_X)^2] = \text{Var}(X)$ and $\mathbb{E}[(Y-\mu_Y)^2] = \text{Var}(Y)$.
		The middle term is the canonical definition of the joint covariance weighted by $2cd$:
		\begin{align*}
			\mathbb{E}\left[ (X - \mu_X)(Y - \mu_Y) \right] \equiv \text{Cov}(X, Y).
		\end{align*}
		Under the strict hypothesis of the statement, the variables $X$ and $Y$ are statistically and functionally independent, which axiomatically implies that the covariance is identically zero ($\text{Cov}(X, Y) = \mathbb{E}[XY] - \mathbb{E}[X]\mathbb{E}[Y] = 0$).
		With the entire intermediate cross term vanishing ($\pm 2cd \cdot 0 = 0$), the formal proof concludes:
		\begin{align*}
			\text{Var}(cX \pm dY) = c^2 \text{Var}(X) + d^2 \text{Var}(Y).
		\end{align*}
		Notice how the intermediate original sign $\pm$ invariably transforms into a positive sum in the variance ($+d^2\text{Var}(Y)$) due to the quadratic scaling of the constants.
	\end{enumerate}
\end{solucion}

\begin{solucion}
	\label{sol:3.3.9}
	For Problem \ref{prob:3.3.9}:
	\begin{enumerate}
		\item The fair-coin game stops at the first heads that appears on the $k$-th toss. The random variable modeling the number of attempts until the first success is $K \sim \text{Geom}(p = 0.5)$, whose probability mass function on the support $\mathcal{S}_K = \set{1, 2, 3, \dots, \infty}$ is:
		\begin{align*}
			P(K = k) = (1 - p)^{k-1} p = (0.5)^{k-1} (0.5) = \left(\frac{1}{2}\right)^k = 2^{-k}.
		\end{align*}
		The monetary payoff awarded by the casino at the end of the experiment is $X = 2^K$. By the Law of the Unconscious Statistician (LOTUS), the mathematical expectation of the reward is computed by summing over the entire countably infinite support:
		\begin{align*}
			\mathbb{E}[X] &= \sum_{k=1}^\infty g(k) \cdot P(K = k) \\
			&= \sum_{k=1}^\infty (2^k) \cdot (2^{-k}) \\
			&= \sum_{k=1}^\infty 2^0 = \sum_{k=1}^\infty 1.
		\end{align*}
		We evaluate the asymptotic behavior of this constant infinite series:
		\begin{align*}
			\mathbb{E}[X] = 1 + 1 + 1 + 1 + \dots = \lim_{n \to \infty} \sum_{k=1}^n 1 = \lim_{n \to \infty} n = \infty.
		\end{align*}
		It is thereby rigorously proven that the expected monetary value of the prize diverges to **infinity ($\mathbb{E}[X] = \infty$)**.
		Economic interpretation (Daniel Bernoulli's Paradox): If a player based their decisions exclusively on the criterion of gross expected value, they should be willing to pledge or pay any imaginable finite monetary amount (even billions of dollars) just to acquire the entry ticket to a game with infinite mathematical expectation. However, in empirical reality no rational human being would pay more than a few coins to play. This deep historical contradiction motivated the modern development of the neoclassical theory of **diminishing marginal utility of money**, in which what is maximized is not the expectation of the net money $\mathbb{E}[X]$, but rather the expectation of its concave utility $\mathbb{E}[u(X)] \ll \infty$.

		\item In a real-world financial environment, no casino has infinite monetary reserves. If the casino imposes a regulatory maximum bankroll cap equivalent to $N$ consecutive tosses, the truncated prize scheme $X_N$ is bounded above by $M = 2^N$:
		\begin{align*}
			X_N = \begin{cases}
				2^k & \text{if the game ends on attempt } k \le N \text{ (probability } 2^{-k}), \\
				2^N & \text{if heads does not appear in the first } N \text{ tosses (probability } P(K > N)).
			\end{cases}
		\end{align*}
		The survival tail probability of a geometric distribution (not having obtained any heads in $N$ consecutive tosses) is exactly the probability of the intersection event of $N$ independent failures:
		\begin{align*}
			P(K > N) = (1 - p)^N = \left(\frac{1}{2}\right)^N = 2^{-N}.
		\end{align*}
		We now formally compute the regularized expected value of the truncated game by summing the contribution of the $N$ possible successful attempts plus the term for the final bankroll cap:
		\begin{align*}
			\mathbb{E}[X_N] &= \left[ \sum_{k=1}^N (2^k) \cdot (2^{-k}) \right] + (2^N) \cdot P(K > N) \\
			&= \left[ \sum_{k=1}^N 1 \right] + (2^N) \cdot (2^{-N}) \\
			&= N + 2^0 = N + 1.
		\end{align*}
		We have derived a most elegant and compact closed-form expression: **the truncated expected value is linear in the maximum number of tosses ($\mathbb{E}[X_N] = N + 1$)**.

		**Numerical evaluation for $N = 25$ tosses:**
		If the casino's maximum bankroll is $N = 25$ (representing an astronomical maximum prize of $2^{25} = 33,554,432$ gold coins), the exact expected value of the game is:
		\begin{align*}
			\mathbb{E}[X_{25}] = 25 + 1 = 26 \text{ monetary units}.
		\end{align*}
		This astonishing logarithmic compression resolves the paradox in practice: even in the face of a multimillion-dollar bankroll of more than 33.5 million dollars, the fair entry price for the game is worth barely **26 coins**.
	\end{enumerate}
\end{solucion}

\begin{solucion}
	\label{sol:3.3.10}
	For Problem \ref{prob:3.3.10}:
	Let $X$ be a discrete random variable with bounded support $\mathcal{S}_X \subseteq [a, b]$, which axiomatically implies that $P(a \le X \le b) = 1$.
	Consider the algebraic product of the two complementary distances from the support to the point $X$: the factor $(X - a)$ and the factor $(b - X)$.
	Since for any observable outcome $\omega \in \Omega$ it holds with probability one that $a \le X(\omega) \le b$, both differences are simultaneously and strictly non-negative over the entire domain:
	\begin{align*}
		X(\omega) - a \ge 0 \quad \text{and} \quad b - X(\omega) \ge 0 \implies (X - a)(b - X) \ge 0.
	\end{align*}
	We algebraically expand the non-negative quadratic product:
	\begin{align*}
		(X - a)(b - X) &= bX - X^2 - ab + aX \\
		&= (a + b)X - X^2 - ab \ge 0.
	\end{align*}
	By the monotonicity axiom of the linear expectation operator (if a random variable $Y \ge 0$ almost surely, then $\mathbb{E}[Y] \ge 0$), we take mathematical expectation on both sides of the inequality and distribute linearly:
	\begin{align*}
		\mathbb{E}\left[ (a + b)X - X^2 - ab \right] \ge 0 \implies (a + b)\mathbb{E}[X] - \mathbb{E}[X^2] - ab \ge 0.
	\end{align*}
	We isolate the raw second-order moment on the right-hand side, preserving the direction of the bound:
	\begin{align*}
		\mathbb{E}[X^2] \le (a + b)\mu - ab, \quad \text{where } \mu \equiv \mathbb{E}[X].
	\end{align*}
	Now, we invoke the König-Huygens computational formula for the central variance:
	\begin{align*}
		\text{Var}(X) = \mathbb{E}[X^2] - \mu^2.
	\end{align*}
	We substitute the upper bound we have just proved for $\mathbb{E}[X^2]$ into the variance formula:
	\begin{align*}
		\text{Var}(X) \le (a + b)\mu - ab - \mu^2.
	\end{align*}
	We reorganize the right-hand side into a canonical quadratic polynomial with respect to the independent variable $\mu$:
	\begin{align*}
		\text{Var}(X) \le -\mu^2 + (a + b)\mu - ab.
	\end{align*}
	To determine the absolute supremum of this concave parabola with respect to $\mu$, we complete the perfect square. We factor out the negative sign from the first two terms:
	\begin{align*}
		-\mu^2 + (a + b)\mu - ab &= -\left[ \mu^2 - (a + b)\mu \right] - ab.
	\end{align*}
	We add and subtract the square of half the linear coefficient $\left(\frac{a+b}{2}\right)^2$ inside the bracket:
	\begin{align*}
		-\left[ \mu^2 - (a + b)\mu + \left(\frac{a+b}{2}\right)^2 - \left(\frac{a+b}{2}\right)^2 \right] - ab &= -\left( \mu - \frac{a+b}{2} \right)^2 + \frac{(a + b)^2}{4} - ab.
	\end{align*}
	We develop the constant fraction independent of $\mu$ with common denominator $4$:
	\begin{align*}
		\frac{(a + b)^2 - 4ab}{4} = \frac{a^2 + 2ab + b^2 - 4ab}{4} = \frac{a^2 - 2ab + b^2}{4} = \frac{(b - a)^2}{4}.
	\end{align*}
	Substituting this identity into our general variance inequality, we obtain the Popoviciu vertex decomposition:
	\begin{align*}
		\text{Var}(X) \le \frac{(b - a)^2}{4} - \left( \mu - \frac{a+b}{2} \right)^2.
	\end{align*}
	Since the subtracted squared term is always non-negative ($\left( \mu - \frac{a+b}{2} \right)^2 \ge 0$), the subtraction can only decrease or preserve the overall magnitude. By eliminating this negative term, the absolute supremum bound is reached:
	\begin{align*}
		\text{Var}(X) \le \frac{(b - a)^2}{4}.
	\end{align*}

	**Exact characterization of the extremal equality distribution:**
	The absolute maximum bound $\text{Var}(X) = \frac{(b-a)^2}{4}$ is achieved with strict mathematical equality if, and only if, two extreme algebraic conditions are simultaneously satisfied:
	1. The completed quadratic term must vanish ($(\mu - \frac{a+b}{2})^2 = 0 \implies \mu = \frac{a+b}{2}$), which requires that the distribution be symmetrically centered at the exact midpoint of the interval $[a, b]$.
	2. The starting inequality $(X - a)(b - X) \ge 0$ must hold with equality almost everywhere ($P((X - a)(b - X) = 0) = 1$), which forces all of the probability mass of $X$ to be located exclusively at the extreme boundary points $x = a$ and $x = b$.

	Therefore, the unique discrete distribution that attains the variance supremum is the **symmetric bimodal Bernoulli distribution concentrated at the endpoints of the interval**, whose probability mass function is:
	\begin{align*}
		P(X = a) = \frac{1}{2} \quad \text{and} \quad P(X = b) = \frac{1}{2}.
	\end{align*}
	With this distribution, the mean is indeed the midpoint $\mu = \frac{a+b}{2}$ and its central variance is half the squared length $\text{Var}(X) = \frac{1}{2}(a - \frac{a+b}{2})^2 + \frac{1}{2}(b - \frac{a+b}{2})^2 = \frac{(b-a)^2}{4}$, confirming the absolute solidity of the theorem.
\end{solucion}
